\documentclass[11pt,a4paper]{article}
\usepackage[T1]{fontenc}
\usepackage[utf8]{inputenc}
\usepackage{lmodern}
\usepackage{amsmath,amssymb,amsthm,mathtools}
\usepackage[margin=1.05in]{geometry}
\usepackage{xcolor}
\usepackage{hyperref}
\usepackage{enumitem}
\usepackage{booktabs}
\hypersetup{colorlinks=true,linkcolor=blue!50!black,citecolor=blue!50!black,urlcolor=blue!50!black}

\newtheorem{theorem}{Theorem}[section]
\newtheorem{lemma}[theorem]{Lemma}
\newtheorem{proposition}[theorem]{Proposition}
\newtheorem{corollary}[theorem]{Corollary}
\newtheorem{citedtheorem}[theorem]{Cited Theorem}
\theoremstyle{definition}
\newtheorem{definition}[theorem]{Definition}
\newtheorem{conjecture}[theorem]{Conjecture}
\newtheorem{hypothesis}[theorem]{Hypothesis}
\theoremstyle{remark}
\newtheorem{remark}[theorem]{Remark}
\newtheorem{convention}[theorem]{Convention}

\DeclareMathOperator{\Tr}{Tr}
\DeclareMathOperator{\spec}{spec}
\DeclareMathOperator{\supp}{supp}
\DeclareMathOperator{\Cov}{Cov}
\DeclareMathOperator{\Var}{Var}
\DeclareMathOperator{\dist}{dist}
\DeclareMathOperator{\osc}{osc}
\DeclareMathOperator{\gap}{gap}
\newcommand{\R}{\mathbb{R}}
\newcommand{\C}{\mathbb{C}}
\newcommand{\Z}{\mathbb{Z}}
\newcommand{\N}{\mathbb{N}}
\newcommand{\E}{\mathbb{E}}
\newcommand{\cH}{\mathcal{H}}
\newcommand{\cA}{\mathcal{A}}
\newcommand{\cG}{\mathcal{G}}
\newcommand{\cL}{\mathcal{L}}
\newcommand{\cP}{\mathcal{P}}
\newcommand{\cE}{\mathcal{E}}
\newcommand{\cR}{\mathcal{R}}
\newcommand{\hT}{\hat{T}}
\newcommand{\hA}{\hat{A}}
\newcommand{\hB}{\hat{B}}
\newcommand{\dU}{\,dU}
\newcommand{\ip}[2]{\langle #1,#2\rangle}
\newcommand{\norm}[1]{\lVert #1\rVert}
\newcommand{\abs}[1]{\lvert #1\rvert}
\newcommand{\TV}{\mathrm{TV}}

\title{Lattice Yang--Mills theory: foundations for the mass gap problem\\
\large Part I of the project ``Solving Yang--Mills and the Mass Gap''}
\author{}
\date{Session 1 (5 November 2026); living document}

\begin{document}
\maketitle

\begin{abstract}
We set up, with complete proofs, the finite-lattice theory on which every later part of this project
rests: Wilson's measure for a compact gauge group, its gauge invariance, reflection positivity, the
transfer matrix and Hamiltonian, the equivalence of the mass gap with exponential clustering in
Euclidean time, an exact reformulation of the gap as the spectral gap of a reversible Markov chain on
the time-slice configuration space (the ground-state chain), the reduction of the continuum problem to
lattice hypotheses, and the existence of a mass gap uniform in the volume at strong coupling. The
strong-coupling result is proved by a self-contained coupling argument (Dobrushin uniqueness) and,
separately, the polymer (closed-surface) representation of the measure is derived and its cluster
expansion is shown to converge for small~$\beta$ via the Kotecký--Preiss criterion. None of the results
is new; the point is to have them in one consistent framework, with the hypotheses that the open
problem requires made explicit (Section~\ref{sec:reduction}). The last section lists what remains open.
\end{abstract}

\tableofcontents

\section{Introduction and conventions}

The Clay problem asks for the construction of a quantum Yang--Mills theory on $\R^4$ with a mass gap.
The only framework in which a non-abelian gauge theory can presently be defined non-perturbatively with
all the relevant structure (positivity, Hamiltonian, gauge invariance) is Wilson's lattice regularisation.
This document develops that framework from scratch, in the generality we need (any compact group $G$ with
a faithful unitary representation, any dimension $d\ge2$, any finite periodic lattice), and isolates the
two statements that are missing: convergence of the continuum limit and uniformity of the gap.

\begin{convention}[Labels]
Every statement is labelled \emph{Theorem/Lemma/Proposition} (complete proof here), \emph{Cited Theorem}
(proof in the literature, reference given), \emph{Conjecture} (no proof), or \emph{Remark} (heuristic
or commentary). Constants are explicit unless stated otherwise; they are not optimised.
\end{convention}

\begin{convention}[Group]
$G$ is a compact Hausdorff topological group (in applications a compact Lie group, e.g.\ $SU(N)$) with
normalised Haar measure $dU$. We fix a faithful continuous unitary representation $\rho:G\to U(N)$ and
write $\Tr U:=\Tr\rho(U)$. For $G=SU(N)$, $\rho$ is the defining representation. Irreducible unitary
representations are denoted $r$, with dimension $d_r$, matrix coefficients $D^r_{ab}$ and character
$\chi_r=\sum_a D^r_{aa}$; $r=0$ is the trivial representation.
\end{convention}

\begin{convention}[Lattice]
Let $d\ge2$ and $L_0,\dots,L_{d-1}\ge 2$ be integers. The lattice is the torus
$\Lambda=\prod_{\mu=0}^{d-1}\Z/L_\mu\Z$ with unit vectors $e_\mu$. Direction $0$ is Euclidean time.
An (oriented) link is a pair $\ell=(x,\mu)$ with $x\in\Lambda$, $0\le\mu<d$, standing for the segment from
$x$ to $x+e_\mu$; $E(\Lambda)$ denotes the set of links, $|E(\Lambda)|=d|\Lambda|$. A plaquette is a pair
$p=(x,\mu\nu)$ with $\mu<\nu$, standing for the unit square with corners $x,x+e_\mu,x+e_\mu+e_\nu,x+e_\nu$;
$\cP(\Lambda)$ is the set of plaquettes, $|\cP(\Lambda)|=\binom d2|\Lambda|$. The four links of $p$ are
$(x,\mu),(x+e_\mu,\nu),(x+e_\nu,\mu),(x,\nu)$; we write $\ell\in p$ or $p\ni\ell$.
\end{convention}

\begin{convention}[Configurations]
A configuration is a map $U:E(\Lambda)\to G$, $\ell\mapsto U_\ell$; we write $U_\mu(x):=U_{(x,\mu)}$ and
$U_{-\mu}(x+e_\mu):=U_\mu(x)^{-1}$ for the reversed link. The configuration space is
$\Omega_\Lambda:=G^{E(\Lambda)}$ with the product Haar measure $dU:=\prod_\ell dU_\ell$. The plaquette variable is
\begin{equation}\label{eq:plaquette}
U_p:=U_\mu(x)\,U_\nu(x+e_\mu)\,U_\mu(x+e_\nu)^{-1}\,U_\nu(x)^{-1},\qquad p=(x,\mu\nu).
\end{equation}
\end{convention}

\section{Compact groups and the Wilson measure}\label{sec:measure}

\subsection{Facts about compact groups}

We use the following standard facts (see e.g.\ Bröcker--tom Dieck, \emph{Representations of Compact Lie
Groups}, Chapters II--III, or Folland, \emph{A Course in Abstract Harmonic Analysis}, Chapter~5).

\begin{citedtheorem}[Haar measure, Peter--Weyl]\label{thm:pw}
\begin{enumerate}[label=(\roman*)]
\item There is a unique regular Borel probability measure $dU$ on $G$ invariant under left and right
translations; it is also invariant under inversion $U\mapsto U^{-1}$.
\item (Schur orthogonality) For irreducible $r,r'$:
$\int \overline{D^r_{ab}(U)}\,D^{r'}_{a'b'}(U)\dU=\delta_{rr'}\delta_{aa'}\delta_{bb'}/d_r$.
\item The functions $\sqrt{d_r}\,D^r_{ab}$ form an orthonormal basis of $L^2(G,dU)$; the span of all matrix
coefficients is dense in $C(G)$ (uniform norm).
\item For irreducible $r,r'$ and $A,B\in G$:
$\int\chi_r(AU)\,\chi_{r'}(U^{-1}B)\dU=\delta_{rr'}\,\chi_r(AB)/d_r$.
\item For irreducible $r$ and $X,Y\in G$: $\int\chi_r(XUYU^{-1})\dU=\chi_r(X)\chi_r(Y)/d_r$
(from (ii), $\int D^r_{ab}(U)D^r_{cd}(U^{-1})\dU=\int D^r_{ab}(U)\overline{D^r_{dc}(U)}\dU=\delta_{ad}\delta_{bc}/d_r$).
\item Every continuous class function $\phi$ (i.e.\ $\phi(gUg^{-1})=\phi(U)$) has an expansion
$\phi=\sum_r c_r\chi_r$ with $c_r=\int\overline{\chi_r}\,\phi\dU$, converging in $L^2$; if $\phi$ is a
finite linear combination of characters, or more generally if $\sum_r|c_r|d_r<\infty$, it converges uniformly.
\end{enumerate}
\end{citedtheorem}

\subsection{The plaquette weight and its character expansion}

\begin{definition}[Plaquette weight]
For $\beta\ge0$ let $w_\beta:G\to(0,\infty)$,
\begin{equation}\label{eq:w}
w_\beta(U):=\exp\Bigl(\frac{\beta}{N}\,\mathrm{Re}\Tr U\Bigr),\qquad
e^{-\beta}\le w_\beta(U)\le e^{\beta}.
\end{equation}
$w_\beta$ is a continuous class function, invariant under $U\mapsto U^{-1}$ (since $\Tr U^{-1}=\overline{\Tr U}$).
Its character coefficients are $c_r(\beta):=\int\overline{\chi_r(U)}\,w_\beta(U)\dU$ and
$c_0(\beta)=\int w_\beta\dU$.
\end{definition}

\begin{lemma}[Positivity of the character coefficients]\label{lem:cr}
For every $\beta>0$ and every irreducible $r$, $c_r(\beta)>0$.
\end{lemma}

\begin{proof}
Write $t=\beta/(2N)$ and $w_\beta(U)=e^{t\Tr U}e^{t\overline{\Tr U}}$. Expanding both exponentials
(absolutely convergent, uniformly on $G$),
\[
w_\beta(U)=\sum_{n,m\ge0}\frac{t^{n+m}}{n!\,m!}\,(\Tr U)^n\,\overline{(\Tr U)^m}
=\sum_{n,m\ge0}\frac{t^{n+m}}{n!\,m!}\,\chi_{\rho^{\otimes n}\otimes\bar\rho^{\otimes m}}(U),
\]
where $\bar\rho$ is the complex-conjugate representation and we used $\chi_{\rho^{\otimes n}}=(\Tr\rho)^n$,
$\chi_{\bar\rho}=\overline{\chi_\rho}$ and multiplicativity of characters under tensor products. Decomposing
$\rho^{\otimes n}\otimes\bar\rho^{\otimes m}=\bigoplus_r m_r(n,m)\,r$ with multiplicities $m_r(n,m)\in\N$
and integrating against $\overline{\chi_r}$ term by term (uniform convergence),
\[
c_r(\beta)=\sum_{n,m\ge0}\frac{t^{n+m}}{n!\,m!}\,m_r(n,m)\ \ge 0 .
\]
For strict positivity it suffices that $m_r(n,m)>0$ for some $(n,m)$. Suppose not: then for all $n,m$,
$\int\overline{D^r_{ab}}\,(\Tr U)^n\overline{(\Tr U)^m}\dU=0$ for all $a,b$ (the matrix coefficients of $r$ are
orthogonal to all matrix coefficients of $\rho^{\otimes n}\otimes\bar\rho^{\otimes m}$, and $(\Tr U)^n\overline{(\Tr U)^m}$
is a sum of such). The algebra $\mathcal{B}$ generated by the functions $\rho(U)_{ij}$, $\overline{\rho(U)_{ij}}$ and $1$
is a $*$-subalgebra of $C(G)$ which separates points because $\rho$ is faithful; by Stone--Weierstrass it is
dense in $C(G)$. Each element of $\mathcal{B}$ is a linear combination of matrix coefficients of the
representations $\rho^{\otimes n}\otimes\bar\rho^{\otimes m}$, hence orthogonal to $D^r_{ab}$; by density,
$D^r_{ab}\perp C(G)$, so $D^r_{ab}=0$, a contradiction.
\end{proof}

\begin{lemma}[The one-link operator]\label{lem:kop}
Let $\beta>0$ and let $k_\beta$ be the integral operator on $L^2(G,dU)$ with kernel $k_\beta(U,V)=w_\beta(UV^{-1})$. Then
\begin{equation}\label{eq:kop}
k_\beta=\sum_r\frac{c_r(\beta)}{d_r}\,\Pi_r ,
\end{equation}
where $\Pi_r$ is the orthogonal projection onto the $d_r^2$-dimensional span of $\{D^r_{ab}\}_{a,b}$ and the
sum converges in Hilbert--Schmidt norm. Consequently $k_\beta$ is self-adjoint, positive and injective, and it is
trace class with
\begin{equation}\label{eq:trace}
\Tr k_\beta=\sum_r c_r(\beta)\,d_r=w_\beta(1)=e^{\beta}.
\end{equation}
\end{lemma}

\begin{proof}
For a unitary representation $D^r(V^{-1})=D^r(V)^*$, so
$\chi_r(UV^{-1})=\sum_{a,b}D^r_{ab}(U)\overline{D^r_{ab}(V)}=\frac1{d_r}\sum_{a,b}e^r_{ab}(U)\overline{e^r_{ab}(V)}$
with the orthonormal functions $e^r_{ab}:=\sqrt{d_r}D^r_{ab}$ of Theorem~\ref{thm:pw}(iii). The class function
$w_\beta$ has the $L^2(G)$-convergent expansion $\sum_rc_r\chi_r$ with $\sum_r|c_r|^2=\norm{w_\beta}_2^2<\infty$
(Parseval), hence $k_\beta(U,V)=\sum_r\frac{c_r}{d_r}\sum_{ab}e^r_{ab}(U)\overline{e^r_{ab}(V)}$ converges in
$L^2(G\times G)$: the terms are orthogonal in $L^2(G\times G)$ and their squared norms sum to
$\sum_r(c_r/d_r)^2d_r^2=\sum_rc_r^2$. This is~\eqref{eq:kop}. Since $c_r>0$ (Lemma~\ref{lem:cr}), $k_\beta$
is positive and injective. The kernel is continuous and the operator is positive, so Mercer's theorem gives
that $k_\beta$ is trace class with $\Tr k_\beta=\int k_\beta(U,U)\dU=w_\beta(1)=e^\beta$; computing the trace from
\eqref{eq:kop} gives $\sum_r(c_r/d_r)\,d_r^2=\sum_rc_rd_r$.
\end{proof}

\begin{corollary}\label{cor:uniform}
The character expansion $w_\beta=\sum_rc_r(\beta)\chi_r$ converges absolutely and uniformly on $G$, and
$1\le c_0(\beta)\le e^\beta$.
\end{corollary}

\begin{proof}
$|c_r\chi_r(U)|\le c_rd_r$ and $\sum_rc_rd_r=e^\beta$ by~\eqref{eq:trace}, so the series converges absolutely
and uniformly to a continuous function which agrees with $w_\beta$ in $L^2$, hence everywhere. Jensen's inequality
and $\int\Tr U\dU=0$ (orthogonality of $\chi_\rho$ to $\chi_0$) give $c_0=\int w_\beta\ge\exp(\frac\beta N\int\mathrm{Re}\Tr U)=1$;
$c_0\le e^\beta$ by~\eqref{eq:w}.
\end{proof}

\begin{remark}
For $G=SU(2)$, $c_j(\beta)=\frac{2(2j+1)}{\beta}I_{2j+1}(\beta)$ ($I_\nu$ = modified Bessel function), so
$c_0=2I_1(\beta)/\beta$ and the fundamental ratio is $u(\beta):=c_{1/2}/(2c_0)=I_2(\beta)/I_1(\beta)=\beta/4+O(\beta^3)$.
These formulas are implemented and tested in \texttt{lattice/ymlattice/exact.py}.
\end{remark}

\subsection{The Wilson measure}

\begin{definition}[Wilson measure]\label{def:wilson}
For $\beta\ge0$ the Wilson weight and partition function are
\[
W_\beta(U):=\prod_{p\in\cP(\Lambda)}w_\beta(U_p),\qquad
Z_\Lambda(\beta):=\int_{\Omega_\Lambda}W_\beta(U)\dU,
\]
and the Wilson measure is the probability measure $d\mu_{\Lambda,\beta}:=Z_\Lambda^{-1}W_\beta\dU$ on $\Omega_\Lambda$.
Expectations are written $\langle F\rangle=\langle F\rangle_{\Lambda,\beta}$. Equivalently
$W_\beta=\exp(-S_\beta)$ with the Wilson action $S_\beta(U)=\beta\sum_p\bigl(1-\frac1N\mathrm{Re}\Tr U_p\bigr)$ up to the
constant $e^{-\beta|\cP|}$, which cancels in $\mu$. With $\beta=2N/g^2$ and $U_\mu(x)=\exp(aA_\mu(x))$ the
action is formally $\frac1{2g^2}\int\Tr F_{\mu\nu}^2\,d^dx+O(a^2)$.
\end{definition}

\begin{theorem}[Basic properties]\label{thm:basic}
\begin{enumerate}[label=(\roman*)]
\item $\mu_{\Lambda,\beta}$ is a well-defined Borel probability measure with continuous density
$W_\beta/Z_\Lambda\in[e^{-2\beta|\cP|},e^{2\beta|\cP|}]$ with respect to $dU$. Thus $\mu_{\Lambda,\beta}$ and $dU$ are mutually
absolutely continuous and every bounded measurable function is integrable.
\item (Gauge invariance) For $g:\Lambda\to G$ define $(g\cdot U)_\mu(x):=g(x)U_\mu(x)g(x+e_\mu)^{-1}$. Then
$(g\cdot U)_p=g(x)U_pg(x)^{-1}$ for $p=(x,\mu\nu)$, hence $W_\beta(g\cdot U)=W_\beta(U)$, and $\mu_{\Lambda,\beta}$ is
invariant under $U\mapsto g\cdot U$.
\item (Elitzur-type vanishing) If $F$ is a bounded measurable function and $g$ a gauge transformation with
$\int_{G^\Lambda}F(g\cdot U)\,dg=0$ for all $U$ (e.g.\ $F=D^r_{ab}(U_\ell)$ for a single link and $r\neq0$),
then $\langle F\rangle=0$. More generally, $\langle F\rangle=\langle\bar F\rangle$ where
$\bar F(U):=\int F(g\cdot U)\,dg$ is the gauge average, so expectations are determined by gauge-invariant
functions.
\item (Lattice symmetries) $\mu_{\Lambda,\beta}$ is invariant under translations of $\Lambda$ and under the
hypercubic group acting on $\Lambda$ and on link orientations (with $U_{-\mu}(x+e_\mu)=U_\mu(x)^{-1}$), and
under the reflections of Section~\ref{sec:rp}.
\end{enumerate}
\end{theorem}

\begin{proof}
(i) $W_\beta$ is continuous and bounded between $e^{-\beta|\cP|}$ and $e^{\beta|\cP|}$ by~\eqref{eq:w}, so
$Z_\Lambda\in[e^{-\beta|\cP|},e^{\beta|\cP|}]$ and the density is as stated.
(ii) Direct computation from~\eqref{eq:plaquette}: the factors $g(x+e_\mu)^{-1}g(x+e_\mu)$, etc., cancel
around the square, leaving $g(x)U_pg(x)^{-1}$; $\mathrm{Re}\Tr$ is conjugation invariant. The map $U\mapsto g\cdot U$
is a product of left and right translations on each factor of $G^{E}$, hence preserves $dU$, and it preserves
$W_\beta$; therefore it preserves $\mu$.
(iii) By (ii), $\langle F\rangle=\langle F\circ g\rangle$ for every $g$; integrating over $g\in G^\Lambda$ with
respect to Haar measure (Fubini) gives $\langle F\rangle=\langle\bar F\rangle$. For $F=D^r_{ab}(U_\ell)$ with
$\ell=(x,\mu)$: $\int D^r_{ab}(g(x)U_\ell g(x+e_\mu)^{-1})\,dg(x)=0$ by Theorem~\ref{thm:pw}(ii) with $r'=0$.
(iv) Both the measure $dU$ and the set of plaquettes are invariant under these maps; under orientation
reversal of a plaquette $U_p\mapsto U_p^{-1}$ (up to conjugation), and $w_\beta(U^{-1})=w_\beta(U)$.
\end{proof}

\begin{remark}[What is and is not an observable]
Statement (iii) is the lattice form of Elitzur's theorem: no gauge-variant local quantity can acquire an
expectation value, at any $\beta$ and in any volume. All physics is in gauge-invariant functions: products
of plaquette variables, Wilson loops $\Tr\prod_{\ell\in\gamma}U_\ell$ along closed paths $\gamma$, and
functions of them. In particular the ``mass gap'' is a property of the gauge-invariant sector only, which is
exactly the sector reconstructed from gauge-invariant correlation functions.
\end{remark}

\section{Reflection positivity}\label{sec:rp}

Throughout this section $L_0=2T$ is even. Let $\theta:\Lambda\to\Lambda$, $\theta(x_0,\vec x)=(-x_0,\vec x)$
(time reflection through the slice $x_0=0$; it also fixes the slice $x_0=T$ pointwise). Put
\[
\Lambda_+:=\{x:0\le x_0\le T\},\qquad \Lambda_-:=\theta\Lambda_+=\{x:T\le x_0\le 2T\},
\]
and $\Lambda_0:=\Lambda_+\cap\Lambda_-=\{x:x_0\in\{0,T\}\}$.
Let $E_\pm$ be the links with both endpoints in $\Lambda_\pm$, $E_0=E_+\cap E_-$ (the spatial links in the two
fixed slices), and similarly $\cP_\pm$ the plaquettes with all corners in $\Lambda_\pm$ and $\cP_0=\cP_+\cap\cP_-$
(the spatial plaquettes in the fixed slices). Every link lies in $E_+\cup E_-$ and every plaquette in
$\cP_+\cup\cP_-$ (a plaquette spans at most one unit of time).

\begin{definition}[Reflection of configurations and functions]
Define the involution $\theta^*:\Omega_\Lambda\to\Omega_\Lambda$ by
\[
(\theta^*U)_{(x,i)}:=U_{(\theta x,i)}\quad(1\le i<d),\qquad
(\theta^*U)_{(x,0)}:=U_{(\theta x-e_0,\,0)}^{-1}.
\]
(The temporal link from $x$ to $x+e_0$ is reflected to the segment from $\theta x$ to $\theta x-e_0$, which is the
reversed link $(\theta x-e_0,0)$.) For a function $F$ on $\Omega_\Lambda$ put $(\Theta F)(U):=\overline{F(\theta^*U)}$.
$\Theta$ is antilinear, involutive and multiplicative. Let $\cA_+$ be the algebra of bounded measurable functions
depending only on $(U_\ell)_{\ell\in E_+}$; then $\Theta\cA_+$ consists of functions of $(U_\ell)_{\ell\in E_-}$.
\end{definition}

\begin{lemma}\label{lem:reflsym}
$\theta^*$ preserves $dU$, and $W_\beta\circ\theta^*=W_\beta$. More precisely, $\theta$ induces a bijection
$p\mapsto\theta p$ of $\cP(\Lambda)$ exchanging $\cP_+$ and $\cP_-$ and fixing $\cP_0$ pointwise, with
$w_\beta((\theta^*U)_p)=w_\beta(U_{\theta p})$ for every $p$.
\end{lemma}

\begin{proof}
$\theta^*$ permutes the link variables and inverts some of them; both operations preserve the product Haar measure
(Theorem~\ref{thm:pw}(i)). For a spatial plaquette $p=(x,ij)$, $(\theta^*U)_p=U_{(\theta x,ij)}$ directly. For a
temporal plaquette $p=(x,0i)$ put $y=\theta x-e_0$, so that $\theta x=y+e_0$ and $\theta(x+e_i)-e_0=y+e_i$; then by
\eqref{eq:plaquette}
\[
(\theta^*U)_p=U_0(y)^{-1}\,U_i(y)\,U_0(y+e_i)\,U_i(y+e_0)^{-1}
=U_0(y)^{-1}\,\bigl(U_{p'}\bigr)^{-1}\,U_0(y),\qquad p'=(y,0i),
\]
since $U_{p'}^{-1}=U_i(y)U_0(y+e_i)U_i(y+e_0)^{-1}U_0(y)^{-1}$. As $w_\beta$ is a class function invariant under
inversion, $w_\beta((\theta^*U)_p)=w_\beta(U_{p'})$. The map $p\mapsto\theta p:=p'$ (resp.\ $(\theta x,ij)$) is a
bijection with the stated properties, so $W_\beta(\theta^*U)=\prod_pw_\beta(U_{\theta p})=W_\beta(U)$.
\end{proof}

\begin{theorem}[Reflection positivity]\label{thm:rp}
For all $F,G\in\cA_+$, the form $(F,G)\mapsto\langle\Theta(F)\,G\rangle_{\Lambda,\beta}$ is sesquilinear and
positive semidefinite; in particular $\langle\Theta(F)F\rangle_{\Lambda,\beta}\ge0$.
\end{theorem}

\begin{proof}
Write $U=(U_0,U_+,U_-)$ with $U_0=(U_\ell)_{\ell\in E_0}$, $U_\pm=(U_\ell)_{\ell\in E_\pm\setminus E_0}$. Then
$W_\beta=W_+W_-/W_0$ with $W_\pm:=\prod_{p\in\cP_\pm}w_\beta(U_p)$, a function of $(U_0,U_\pm)$ only, and
$W_0:=\prod_{p\in\cP_0}w_\beta(U_p)>0$, a function of $U_0$ only. Since $\theta$ is the identity on the fixed
slices, $(\theta^*U)_\ell=U_\ell$ for $\ell\in E_0$, while $\ell\mapsto\theta\ell$ is a bijection
$E_+\setminus E_0\to E_-\setminus E_0$; thus for fixed $U_0$ the map $\Phi_{U_0}:U_-\mapsto\tilde U_+:=((\theta^*U)_\ell)_{\ell\in E_+\setminus E_0}$
is a measure-preserving bijection and $G(\theta^*U)=G(U_0,\Phi_{U_0}(U_-))$ for $G\in\cA_+$. By
Lemma~\ref{lem:reflsym}, $W_-(U_0,U_-)=\prod_{p\in\cP_-}w_\beta(U_p)=\prod_{p\in\cP_+}w_\beta((\theta^*U)_p)=W_+(U_0,\Phi_{U_0}(U_-))$.
Hence, by Fubini and the change of variables $\tilde U_+=\Phi_{U_0}(U_-)$,
\begin{align*}
Z_\Lambda\langle\Theta(F)G\rangle
&=\int dU_0\,W_0(U_0)^{-1}\Bigl[\int dU_+\,G(U_0,U_+)W_+(U_0,U_+)\Bigr]\\
&\qquad\qquad\times\Bigl[\int dU_-\,\overline{F(U_0,\Phi_{U_0}(U_-))}\,W_+(U_0,\Phi_{U_0}(U_-))\Bigr]\\
&=\int dU_0\,W_0(U_0)^{-1}\,\overline{\Psi_F(U_0)}\,\Psi_G(U_0),\qquad
\Psi_F(U_0):=\int dU_+\,F(U_0,U_+)\,W_+(U_0,U_+).
\end{align*}
The right-hand side is manifestly a positive semidefinite sesquilinear form in $(F,G)$.
\end{proof}

\begin{remark}[Other reflections]
(i) By the hypercubic symmetry of Theorem~\ref{thm:basic}(iv), the same statement holds for the reflection
through a plane of sites orthogonal to any axis. In particular, for any real $F\in\cA_+$ that is a function of
links in a slab, $\langle\Theta(F)F\rangle-\langle F\rangle\langle\Theta F\rangle\ge0$ follows from applying
Theorem~\ref{thm:rp} to $F-\langle F\rangle$ (note $\langle\Theta F\rangle=\overline{\langle F\rangle}$ by
Lemma~\ref{lem:reflsym}). This is the lattice statement that ``reflection-symmetric'' two-point functions are
non-negative, used in the numerical checks. (ii) Reflection positivity also holds for reflections through planes
bisecting links (Osterwalder--Seiler); it is not needed here, because the transfer matrix of
Section~\ref{sec:transfer} lives on slices of sites.
\end{remark}

\section{Transfer matrix, Hamiltonian, mass gap}\label{sec:transfer}

\subsection{The time-slice Hilbert space}

Let $\Lambda_s:=\prod_{i=1}^{d-1}\Z/L_i\Z$ be the spatial lattice, $E_s$ its (spatial) links, $\cP_s$ its
plaquettes, $\Omega_s:=G^{E_s}$ with product Haar measure $dU$, and
\[
\cH:=L^2(\Omega_s,dU).
\]
The spatial gauge group $\cG_s:=G^{\Lambda_s}$ acts on $\Omega_s$ by $(g\cdot U)_{(x,i)}=g(x)U_{(x,i)}g(x+e_i)^{-1}$
and on $\cH$ by the unitary operators $(R_g\psi)(U):=\psi(g^{-1}\cdot U)$; $g\mapsto R_g$ is a strongly continuous
unitary representation. The \emph{physical} (gauge-invariant) subspace is
\[
\cH_{\rm phys}:=\{\psi\in\cH:\ R_g\psi=\psi\ \ \forall g\in\cG_s\},\qquad
P:=\int_{\cG_s}R_g\,dg ,
\]
where the integral (over normalised Haar measure of the compact group $\cG_s$) converges strongly. $P$ is the
orthogonal projection onto $\cH_{\rm phys}$: $P^2=P$ by invariance of Haar measure, $P^*=P$ by inversion invariance,
and $R_hP=P$ for all $h$ gives $\mathrm{ran}\,P=\cH_{\rm phys}$.

\subsection{Definition of the transfer matrix}

\begin{definition}\label{def:T}
Let $\beta>0$. Define the following operators on $\cH$.
\begin{enumerate}[label=(\roman*)]
\item $V$: multiplication by $V(U):=\prod_{p\in\cP_s}w_\beta(U_p)^{1/2}$ (a continuous function with values in
$[e^{-\beta|\cP_s|/2},e^{\beta|\cP_s|/2}]$). $V$ is bounded, self-adjoint, strictly positive and commutes with every $R_g$.
\item $K_1:=\bigotimes_{\ell\in E_s}k_\beta$, the integral operator with kernel
$K_1(U,U')=\prod_{\ell\in E_s}w_\beta(U_\ell U_\ell'^{-1})$.
\item $K:=K_1P$.
\item The \emph{transfer matrix} $\hT:=VKV$.
\end{enumerate}
\end{definition}

\begin{lemma}\label{lem:K}
\begin{enumerate}[label=(\roman*)]
\item $K_1$ is self-adjoint, positive, injective and trace class with $\Tr K_1=e^{\beta|E_s|}$; it commutes with every $R_g$.
\item $K=K_1P=PK_1=PK_1P$, and $K$ is the integral operator with the continuous kernel
\begin{equation}\label{eq:Kkernel}
K(U,U')=\int_{\cG_s}dg\,K_1(U,g\cdot U')=\int\prod_{x\in\Lambda_s}dV_x\prod_{(x,i)\in E_s}
w_\beta\bigl(U_{(x,i)}\,V_{x+e_i}\,U'^{-1}_{(x,i)}\,V_x^{-1}\bigr)\ \ge e^{-\beta|E_s|}.
\end{equation}
\item $\hT=P\,(VK_1V)\,P$ is self-adjoint, positive and trace class, $\hT=P\hT P$, and $\hT$ is injective on
$\cH_{\rm phys}$. Its kernel $T(U,U'):=V(U)K(U,U')V(U')$ is continuous and strictly positive on $\Omega_s\times\Omega_s$.
\end{enumerate}
\end{lemma}

\begin{proof}
(i) $K_1$ is a tensor product of copies of $k_\beta$ (Lemma~\ref{lem:kop}); tensor products of positive injective
trace class operators are positive, injective and trace class, with $\Tr K_1=(\Tr k_\beta)^{|E_s|}=e^{\beta|E_s|}$.
Gauge covariance of the kernel: $(g\cdot U)_\ell(g\cdot U')_\ell^{-1}=g(x)U_\ell U_\ell'^{-1}g(x)^{-1}$ for
$\ell=(x,i)$, and $w_\beta$ is a class function; hence $K_1(g\cdot U,g\cdot U')=K_1(U,U')$, i.e.\ $R_gK_1R_g^{-1}=K_1$.
(ii) By (i), $K_1$ commutes with $P$, so $K_1P=PK_1=PK_1P$. For the kernel: for $\psi\in\cH$,
$(K_1P\psi)(U)=\int dU'K_1(U,U')\int dg\,\psi(g^{-1}\cdot U')=\int dg\int dU''K_1(U,g\cdot U'')\psi(U'')$ by the
substitution $U'=g\cdot U''$ (measure preserving), which is \eqref{eq:Kkernel} with $V_x=g(x)$ after writing out
$(g\cdot U')^{-1}_{(x,i)}=g(x+e_i)U'^{-1}_{(x,i)}g(x)^{-1}$; the integrand is continuous and bounded below by $e^{-\beta|E_s|}$,
and so is the integral. Continuity of $K$ follows by dominated convergence.
(iii) $V$ commutes with $P$ (spatial plaquettes are gauge invariant), hence $\hT=VPK_1PV=P(VK_1V)P$. $VK_1V$ is
self-adjoint, positive and trace class (two-sided ideal), hence so is $\hT$, and $\hT=P\hT P$. If $\psi\in\cH_{\rm phys}$
then $\ip\psi{\hT\psi}=\ip{V\psi}{K_1V\psi}>0$ unless $V\psi=0$, i.e.\ $\psi=0$; so $\hT$ is injective on
$\cH_{\rm phys}$. The kernel statement is immediate from (ii).
\end{proof}

\subsection{Partition function and correlation functions as traces}

For a bounded measurable $A:\Omega_s\to\C$ let $\hA$ denote multiplication by $A$ on $\cH$, and for $0\le t<L_0$
let $A(t)$ denote the function on $\Omega_\Lambda$ obtained by evaluating $A$ on the spatial links of the time slice
$x_0=t$.

\begin{theorem}[Transfer matrix formalism]\label{thm:tm}
Let $L_0\ge2$. Then $Z_\Lambda(\beta)=\Tr\hT^{L_0}$, and for bounded measurable $A_1,\dots,A_n$ on $\Omega_s$ and
$0\le t_1<t_2<\dots<t_n\le L_0-1$,
\begin{equation}\label{eq:traces}
\Bigl\langle\prod_{k=1}^nA_k(t_k)\Bigr\rangle_{\Lambda,\beta}
=\frac{\Tr\bigl(\hT^{\,L_0-t_n}\hA_n\hT^{\,t_n-t_{n-1}}\hA_{n-1}\cdots\hA_2\hT^{\,t_2-t_1}\hA_1\hT^{\,t_1}\bigr)}{\Tr\hT^{L_0}} .
\end{equation}
\end{theorem}

\begin{proof}
Label the spatial links of slice $t$ by $U^{(t)}\in\Omega_s$ and the temporal links from slice $t$ to $t+1$ by
$V^{(t)}\in G^{\Lambda_s}$; then $\Omega_\Lambda=\prod_t(\Omega_s\times G^{\Lambda_s})$ and $dU$ is the product measure.
The Wilson weight factorises as
\[
W_\beta(U)=\prod_{t}\Bigl[\prod_{p\in\cP_s}w_\beta(U^{(t)}_p)\Bigr]\prod_t\Bigl[\prod_{(x,i)\in E_s}
w_\beta\bigl(U^{(t)}_{(x,i)}V^{(t)}_{x+e_i}(U^{(t+1)}_{(x,i)})^{-1}(V^{(t)}_x)^{-1}\bigr)\Bigr],
\]
where the first bracket is $V(U^{(t)})^2$ and the second is the weight of the temporal plaquettes between slices $t$
and $t+1$ (the plaquette $(x,0i)$ at time $t$ has $U_p=V_xU'_{(x,i)}V_{x+e_i}^{-1}U_{(x,i)}^{-1}$, whose inverse is
conjugate to the argument written, and $w_\beta$ is a class function invariant under inversion). Integrating over all
$V^{(t)}$ produces $\prod_tK(U^{(t)},U^{(t+1)})$ by \eqref{eq:Kkernel}, so
\[
\int\prod_kA_k(U^{(t_k)})\,W_\beta\dU
=\int\prod_{t}dU^{(t)}\ \prod_kA_k(U^{(t_k)})\ \prod_{t=0}^{L_0-1}T(U^{(t)},U^{(t+1)}),
\]
with $U^{(L_0)}\equiv U^{(0)}$.
For Hilbert--Schmidt operators $X,Y$ with kernels $X(U,U'),Y(U,U')$ one has $\Tr(XY)=\iint X(U,U')Y(U',U)\dU\dU'$, and the
kernel of a product of integral operators is the composition of kernels. Since $L_0\ge2$, the operator inside the trace
in \eqref{eq:traces} can be written as a product $XY$ of two factors each containing at least one factor $\hT$ (which
is Hilbert--Schmidt, being trace class) and bounded multiplication operators (with kernels $A(U)\delta(U-U')$, which act
on kernels by pointwise multiplication). Applying the trace formula and expanding the compositions yields exactly the
integral above. With no insertions this gives $Z_\Lambda=\Tr\hT^{L_0}$.
\end{proof}

\subsection{Perron--Frobenius, the vacuum, and the Hamiltonian}

\begin{theorem}[Vacuum]\label{thm:vacuum}
Let $\beta>0$. The operator norm $\lambda_0:=\norm{\hT}$ is a simple eigenvalue of $\hT$, and its eigenvector can be
chosen as a continuous function $\Omega$ with $\Omega(U)>0$ for all $U\in\Omega_s$ and $\norm\Omega=1$. Moreover
$\Omega\in\cH_{\rm phys}$, and every other eigenvalue of $\hT$ is strictly smaller than $\lambda_0$.
\end{theorem}

\begin{proof}
$\hT$ is compact, self-adjoint and positive, so $\lambda_0=\norm{\hT}=\max\spec\hT$ is an eigenvalue, and
$\ip\varphi{\hT\varphi}=\lambda_0\norm\varphi^2$ holds iff $\hT\varphi=\lambda_0\varphi$ (because $\lambda_0-\hT\ge0$
and $\ip\varphi{(\lambda_0-\hT)\varphi}=0$ forces $(\lambda_0-\hT)^{1/2}\varphi=0$).

\emph{Positivity improving.} By Lemma~\ref{lem:K}(iii), $T(U,U')\ge c>0$ for all $U,U'$. Hence if $\psi\ge0$ and $\psi\neq0$
in $\cH$, then $(\hT\psi)(U)=\int T(U,U')\psi(U')\dU'\ge c\int\psi>0$ for every $U$.

\emph{Eigenvectors are positive up to a phase.} Let $\hT\psi=\lambda_0\psi$, $\psi\neq0$. Since $T>0$,
\begin{align*}
\lambda_0\norm\psi^2=\ip\psi{\hT\psi}&=\mathrm{Re}\iint\overline{\psi(U)}T(U,U')\psi(U')
\le\iint|\psi(U)|T(U,U')|\psi(U')|\\
&=\ip{|\psi|}{\hT|\psi|}\le\lambda_0\norm{|\psi|}^2=\lambda_0\norm\psi^2 ,
\end{align*}
so all inequalities are equalities. The last one gives $\hT|\psi|=\lambda_0|\psi|$, hence $|\psi|=\lambda_0^{-1}\hT|\psi|>0$
everywhere by positivity improving. The first one gives $\iint T(U,U')\bigl[|\psi(U)||\psi(U')|-\mathrm{Re}(\overline{\psi(U)}\psi(U'))\bigr]=0$
with non-negative integrand, so $\overline{\psi(U)}\psi(U')=|\psi(U)||\psi(U')|$ for a.e.\ $(U,U')$, i.e.\ the phase of
$\psi$ is a.e.\ constant: $\psi=e^{i\alpha}|\psi|$.

\emph{Simplicity.} If $\psi_1,\psi_2$ are eigenvectors for $\lambda_0$, we may take both strictly positive. Then
$\psi_1-c\psi_2$ with $c=\ip{\psi_2}{\psi_1}/\norm{\psi_2}^2$ is an eigenvector orthogonal to $\psi_2$; if it were non-zero
it would be a phase times a strictly positive function, which cannot be orthogonal to $\psi_2>0$. So $\psi_1=c\psi_2$.

\emph{Continuity and gauge invariance.} $\Omega=\lambda_0^{-1}\hT\Omega$ and $T$ is continuous on the compact space
$\Omega_s\times\Omega_s$, so $\Omega$ is continuous by dominated convergence. Since $\hT=P\hT P$,
$\Omega=\lambda_0^{-1}P\hT P\Omega\in\mathrm{ran}\,P=\cH_{\rm phys}$. Finally the eigenvalue $\lambda_0$ of the
self-adjoint compact operator $\hT$ is isolated and simple, so all other eigenvalues are strictly smaller.
\end{proof}

\begin{definition}[Hamiltonian and mass gap]\label{def:H}
Let $\lambda_0>\lambda_1\ge\lambda_2\ge\cdots>0$ be the eigenvalues of $\hT|_{\cH_{\rm phys}}$ (counted with
multiplicity; they are positive by Lemma~\ref{lem:K}(iii) and accumulate only at $0$ by compactness). The
\emph{Hamiltonian} is the self-adjoint operator
\[
H:=-\log\bigl(\hT/\lambda_0\bigr)\quad\text{on }\cH_{\rm phys},\qquad \spec H=\{E_n:=\log(\lambda_0/\lambda_n)\},
\]
so that $0=E_0<E_1\le E_2\le\cdots$, $E_n\to\infty$, $H\Omega=0$, and $e^{-tH}=(\hT/\lambda_0)^t$ for $t\in\N$. The
\emph{mass gap} of the finite lattice $\Lambda_s$ at coupling $\beta$ is
\[
m_{\Lambda_s}(\beta):=E_1=\log\frac{\lambda_0}{\lambda_1}>0 .
\]
(In physical units, with lattice spacing $a$, the Hamiltonian is $H/a$ and the gap $m/a$.)
\end{definition}

\begin{remark}
The strict positivity $m_{\Lambda_s}(\beta)>0$ is a statement about one finite lattice. It carries no information
about $\inf_{\Lambda_s}m_{\Lambda_s}(\beta)$, which is what the mass gap problem is about; see
Section~\ref{sec:strong} for the one regime where uniformity is known.
\end{remark}

\subsection{Infinite temporal extent}

\begin{theorem}[Ground-state expectations]\label{thm:gs}
Let $A,B$ be bounded measurable gauge-invariant functions on $\Omega_s$ and $t\in\N$. Then, keeping $\Lambda_s$ and
$\beta$ fixed and letting $L_0\to\infty$ (with $t<L_0$),
\[
\lim_{L_0\to\infty}\langle A(t)B(0)\rangle_{\Lambda,\beta}=\ip\Omega{\hA e^{-tH}\hB\Omega},\qquad
\lim_{L_0\to\infty}\langle A(0)\rangle_{\Lambda,\beta}=\ip\Omega{\hA\Omega}=\int A\,\Omega^2\dU .
\]
\end{theorem}

\begin{proof}
Gauge invariance of $A$ gives $R_g\hA R_g^{-1}=\hA$, hence $\hA$ commutes with $P$ and maps $\cH_{\rm phys}$ into
itself; so $\hB\Omega\in\cH_{\rm phys}$ and $e^{-tH}\hB\Omega=(\hT/\lambda_0)^t\hB\Omega$ is well defined. By
Theorem~\ref{thm:tm},
$\langle A(t)B(0)\rangle=\Tr\bigl((\hT/\lambda_0)^{L_0-t}\hA(\hT/\lambda_0)^t\hB\bigr)/\Tr(\hT/\lambda_0)^{L_0}$.
Write $(\hT/\lambda_0)^n=|\Omega\rangle\langle\Omega|+R_n$ on $\cH$ (spectral theorem, Theorem~\ref{thm:vacuum});
$R_n$ is trace class with $\norm{R_n}_1=\sum_{k\ge1}(\lambda_k/\lambda_0)^n\to0$ as $n\to\infty$ by dominated
convergence (the terms are dominated by $\lambda_k/\lambda_0$, which are summable since $\hT$ is trace class).
Using $|\Tr(RX)|\le\norm R_1\norm X$ for bounded $X$, the numerator converges to
$\Tr(|\Omega\rangle\langle\Omega|\hA(\hT/\lambda_0)^t\hB)=\ip\Omega{\hA(\hT/\lambda_0)^t\hB\Omega}$ and the
denominator to $1$.
\end{proof}

\subsection{Mass gap and exponential clustering}

Write $\nu(dU):=\Omega(U)^2\dU$ (a probability measure on $\Omega_s$, the \emph{vacuum measure}), $\langle A\rangle_\nu:=\int A\,d\nu$.

\begin{theorem}[Gap $\Leftrightarrow$ clustering in time]\label{thm:gapclust}
Fix $\Lambda_s$ and $\beta>0$.
\begin{enumerate}[label=(\roman*)]
\item For all bounded measurable gauge-invariant $A,B$ and all $t\in\N$,
\[
\bigl|\ip\Omega{\hA e^{-tH}\hB\Omega}-\langle A\rangle_\nu\langle B\rangle_\nu\bigr|
\le e^{-m_{\Lambda_s}(\beta)\,t}\,\Var_\nu(A)^{1/2}\Var_\nu(B)^{1/2}
\le e^{-m_{\Lambda_s}(\beta)\,t}\,\norm A_\infty\norm B_\infty .
\]
\item Conversely, suppose $m>0$ and for every bounded measurable gauge-invariant $A$ there is $C_A<\infty$ with
\[
\bigl|\ip\Omega{\hA^*e^{-tH}\hA\Omega}-|\langle A\rangle_\nu|^2\bigr|\le C_A\,e^{-mt}\qquad\text{for all }t\in\N .
\]
Then $m_{\Lambda_s}(\beta)\ge m$.
\end{enumerate}
\end{theorem}

\begin{proof}
(i) Let $Q:=1-|\Omega\rangle\langle\Omega|$ on $\cH_{\rm phys}$. Then $e^{-tH}=|\Omega\rangle\langle\Omega|+e^{-tH}Q$
and $\norm{e^{-tH}Q}=e^{-E_1t}$. Since $\ip\Omega{\hA e^{-tH}\hB\Omega}=\ip{\hA^*\Omega}{e^{-tH}\hB\Omega}$ and
$\ip{\hA^*\Omega}{\Omega}\ip\Omega{\hB\Omega}=\langle A\rangle_\nu\langle B\rangle_\nu$, the connected expectation
equals $\ip{Q\hA^*\Omega}{e^{-tH}Q\hB\Omega}$, bounded by $e^{-E_1t}\norm{Q\hA^*\Omega}\norm{Q\hB\Omega}$, and
$\norm{Q\hB\Omega}^2=\int|B-\langle B\rangle_\nu|^2\,d\nu=\Var_\nu(B)\le\norm B_\infty^2$.

(ii) Let $\psi:=Q\hA\Omega$ and let $\mu_\psi$ be the spectral measure of $H$ for $\psi$; it is supported in
$[E_1,\infty)$, and $\ip\psi{e^{-tH}\psi}=\int e^{-tE}\,d\mu_\psi(E)$ is the connected expectation in the hypothesis.
For $0<\epsilon<m$,
\[
\mu_\psi\bigl([0,m-\epsilon]\bigr)\,e^{-t(m-\epsilon)}\le\int e^{-tE}\,d\mu_\psi(E)\le C_Ae^{-mt},
\]
so $\mu_\psi([0,m-\epsilon])\le C_Ae^{-\epsilon t}\to0$: $\psi$ is orthogonal to the spectral subspace of $H$ for
$[0,m-\epsilon]$. Suppose $E_1<m$ and pick $\epsilon$ with $E_1<m-\epsilon$; then every $\psi=Q\hA\Omega$ is orthogonal
to the (non-zero) eigenspace $\cH_1$ of $E_1$. Now $\{\hA\Omega: A\text{ bounded, gauge invariant}\}$ is dense in
$\cH_{\rm phys}$: bounded gauge-invariant functions are dense in $\cH_{\rm phys}$ (truncation $\psi\mapsto\psi\mathbf 1_{|\psi|\le M}$
preserves gauge invariance), and for bounded gauge-invariant $\varphi$ the function $A:=\varphi/\Omega$ is bounded
(as $\min\Omega>0$ by continuity and compactness) and gauge invariant with $\hA\Omega=\varphi$. Since $Q$ is continuous and
surjective onto $\Omega^\perp\cap\cH_{\rm phys}\supseteq\cH_1$, the vectors $Q\hA\Omega$ are dense in $\Omega^\perp$,
so $\cH_1\perp\Omega^\perp$, forcing $\cH_1=0$, a contradiction. Hence $E_1\ge m$.
\end{proof}

\begin{corollary}[Gap from Euclidean correlations]\label{cor:gapfromcorr}
Suppose $m>0$ and $C<\infty$ are such that for all $L_0$, all $t$ with $0\le t\le L_0/2$ and all bounded gauge-invariant
$A$ on $\Omega_s$,
\[
\bigl|\langle\overline{A}(t)A(0)\rangle_{\Lambda,\beta}-\langle\overline{A}(t)\rangle_{\Lambda,\beta}\langle A(0)\rangle_{\Lambda,\beta}\bigr|
\le C_A\,e^{-mt}
\]
with $C_A$ independent of $L_0$. Then $m_{\Lambda_s}(\beta)\ge m$.
\end{corollary}

\begin{proof}
Let $L_0\to\infty$ at fixed $t$ in Theorem~\ref{thm:gs} (applied to $\bar A$, $A$ and to $A$, $1$) and apply
Theorem~\ref{thm:gapclust}(ii).
\end{proof}

\subsection{The ground-state Markov chain}\label{sec:chain}

\begin{theorem}[Exact reformulation of the gap]\label{thm:chain}
Define $P(U,U'):=\dfrac{T(U,U')\,\Omega(U')}{\lambda_0\,\Omega(U)}$ on $\Omega_s\times\Omega_s$. Then:
\begin{enumerate}[label=(\roman*)]
\item $P$ is a continuous, strictly positive Markov kernel: $\int P(U,U')\dU'=1$ for every $U$.
\item $P$ is reversible with respect to the vacuum measure $\nu=\Omega^2dU$: $\Omega(U)^2P(U,U')=\Omega(U')^2P(U',U)$.
\item The Markov operator $(P\varphi)(U)=\int P(U,U')\varphi(U')\dU'$ on $L^2(\nu)$ is unitarily equivalent to
$\hT/\lambda_0$ on $L^2(dU)$ via $M:L^2(\nu)\to L^2(dU)$, $M\varphi=\Omega\varphi$. Consequently $P$ is self-adjoint on
$L^2(\nu)$, $1$ is a simple eigenvalue with eigenfunction $\mathbf 1$, and with
\[
\gap(P):=1-\sup\Bigl\{\ip\varphi{P\varphi}_\nu:\ \norm\varphi_\nu=1,\ \int\varphi\,d\nu=0\Bigr\}
=\inf_{\varphi\ne\text{const}}\frac{\ip\varphi{(1-P)\varphi}_\nu}{\Var_\nu(\varphi)}
\]
one has
\begin{equation}\label{eq:gapchain}
m_{\Lambda_s}(\beta)=-\log\bigl(1-\gap(P)\bigr),\qquad \gap(P)=1-\lambda_1/\lambda_0 .
\end{equation}
\end{enumerate}
\end{theorem}

\begin{proof}
(i) $\int P(U,U')\dU'=(\hT\Omega)(U)/(\lambda_0\Omega(U))=1$; continuity and positivity are inherited from $T$ and $\Omega$.
(ii) Both sides equal $\Omega(U)T(U,U')\Omega(U')/\lambda_0$, which is symmetric since $T(U,U')=T(U',U)$ ($\hT$ is
self-adjoint with a real continuous kernel). (iii) $M$ is unitary because $\norm{M\varphi}^2_{L^2(dU)}=\int\Omega^2|\varphi|^2=\norm\varphi_\nu^2$,
and $(MPM^{-1}\psi)(U)=\Omega(U)\int P(U,U')\psi(U')/\Omega(U')\dU'=\lambda_0^{-1}\int T(U,U')\psi(U')\dU'$. Hence
$\spec P=\spec(\hT)/\lambda_0=\{1,\lambda_1/\lambda_0,\dots\}\cup\{0\}$ with $1$ simple (Theorem~\ref{thm:vacuum}),
$M^{-1}\Omega=\mathbf 1$, and the supremum in the definition of $\gap(P)$ is the second largest eigenvalue
$\lambda_1/\lambda_0$ (the eigenvalue $0$ on $\cH_{\rm phys}^\perp$ being smaller). The variational formula and
\eqref{eq:gapchain} follow from Definition~\ref{def:H}.
\end{proof}

\begin{remark}[Structure of the chain]\label{rem:chain}
Unfolding the definitions,
\[
P(U,U')=\frac{1}{\lambda_0}\,\frac{\Omega(U')}{\Omega(U)}\prod_{p\in\cP_s}w_\beta(U_p)^{1/2}w_\beta(U'_p)^{1/2}
\int\prod_xdV_x\prod_{(x,i)}w_\beta\bigl(U_{(x,i)}V_{x+e_i}U'^{-1}_{(x,i)}V_x^{-1}\bigr).
\]
Every factor except the ratio $\Omega(U')/\Omega(U)$ is a product of local terms: the chain is a gauge-averaged
parallel update of all spatial links, reweighted by the vacuum wave functional. All the non-locality, and hence
the whole difficulty of the uniform gap, sits in $\Omega$. For spin systems, uniform spectral gaps of such local
reversible dynamics are equivalent to strong mixing properties of the stationary measure (Stroock--Zegarlinski;
Martinelli--Olivieri). The mass gap problem is, in this exact sense, a question about the mixing properties of the
$(d-1)$-dimensional vacuum measure $\nu=\Omega^2dU$, uniformly in the volume. This reformulation is used in
\texttt{docs/03-research-program.md}, \S3.5.
\end{remark}

\section{From lattice hypotheses to the continuum theorem}\label{sec:reduction}

This section makes precise the (folklore) statement that the Clay problem reduces to two properties of the
lattice theory: convergence of correlation functions and uniform exponential clustering. Nothing here is hard;
the value is in the explicit list of hypotheses.

\subsection{Setting}

Fix $d=4$ and a sequence of lattice spacings $a_n\downarrow0$, couplings $\beta_n>0$ and periodic lattices
$\Lambda_n$ with $L^{(n)}_\mu a_n\to\infty$ for every $\mu$ (we also allow $L^{(n)}_\mu a_n=L$ fixed in a first step).
Points of $a_n\Lambda_n$ are identified with points of the torus $\prod_\mu(\R/L^{(n)}_\mu a_n\Z)$ and, for
compactly supported test functions, with points of $\R^4$ for $n$ large.

A \emph{family of local observables} is a set of indices $\alpha\in I$ and, for each $n$, $\alpha$ and
$x\in\Lambda_n$, a bounded measurable gauge-invariant function $\Phi^{(n)}_\alpha(x)$ on $\Omega_{\Lambda_n}$ depending only on
links within a fixed lattice distance $R_\alpha$ of $x$, covariant under lattice translations, and such that
$\Phi^{(n)}_\alpha(x)$ is real and \emph{reflection covariant}: $\Theta\Phi^{(n)}_\alpha(x)=\Phi^{(n)}_\alpha(\theta x)$ for
the reflections $\Theta$ of Section~\ref{sec:rp} (for every fixed slice). Examples: suitably renormalised plaquette
sums $\Phi^{(n)}(x)=Z_n\bigl(\sum_{p\ni x}(1-\tfrac1N\mathrm{Re}\Tr U_p)-\langle\cdot\rangle\bigr)/a_n^4$, which should
converge to $\Tr F_{\mu\nu}F_{\mu\nu}(x)$, and Wilson loops of fixed shape in lattice units. The lattice Schwinger
functions are
\begin{gather*}
S^{(n)}_{\alpha_1\cdots\alpha_k}(x_1,\dots,x_k):=\Bigl\langle\prod_{i=1}^k\Phi^{(n)}_{\alpha_i}(x_i)\Bigr\rangle_{\Lambda_n,\beta_n},\\
S^{(n)}(f_1\otimes\cdots\otimes f_k):=\sum_{x_1,\dots,x_k\in a_n\Lambda_n}a_n^{4k}\prod_if_i(x_i)\,S^{(n)}(x_1,\dots,x_k)
\end{gather*}
for test functions $f_i\in\mathcal S(\R^4)$ (compactly supported, for $n$ large enough).

\subsection{Hypotheses}

\begin{hypothesis}[L1: convergence and Euclidean invariance]\label{hyp:L1}
For every $k$ and $\alpha_1,\dots,\alpha_k$, and all $f_1,\dots,f_k\in C_c^\infty(\R^4)$, the limit
$S_{\alpha_1\cdots\alpha_k}(f_1\otimes\cdots\otimes f_k):=\lim_nS^{(n)}_{\alpha_1\cdots\alpha_k}(f_1\otimes\cdots\otimes f_k)$
exists, extends to a tempered distribution on $\R^{4k}$ satisfying the Osterwalder--Schrader regularity axiom (OS0)
including the linear growth condition, and is invariant under the full Euclidean group $E(4)$ (OS1).
\end{hypothesis}

\begin{hypothesis}[L2: uniform exponential clustering]\label{hyp:L2}
There exists $\Delta>0$ such that for all $k,l$, all indices, and all $f_i,g_j\in C_c^\infty(\R^4)$ there is a constant
$C<\infty$ (depending on the $f_i,g_j$ but not on $n$) with
\[
\bigl|S^{(n)}(f_1\otimes\cdots\otimes f_k\otimes g_{1,t}\otimes\cdots\otimes g_{l,t})
-S^{(n)}(f_1\otimes\cdots\otimes f_k)\,S^{(n)}(g_1\otimes\cdots\otimes g_l)\bigr|\le C\,e^{-\Delta t}
\]
for all $n$ and all $t\in a_n\N$ with $0\le t\le L^{(n)}_0a_n/2$, where $g_{j,t}(x):=g_j(x-te_0)$.
\end{hypothesis}

\begin{hypothesis}[L3: non-triviality]\label{hyp:L3}
Some truncated (connected) four-point function of the limit $S$ is not identically zero.
\end{hypothesis}

\subsection{The reduction theorem}

\begin{citedtheorem}[Osterwalder--Schrader reconstruction]\label{thm:OS}
A family of Schwinger functions satisfying (OS0) with the linear growth condition, (OS1) Euclidean invariance,
(OS2) reflection positivity, (OS3) symmetry and (OS4) clustering determines a Hilbert space $\cH$, a unique (up to
phase) vacuum vector $\Omega$, a strongly continuous unitary positive-energy representation of the Poincaré group with
generator of time translations $H\ge0$, $H\Omega=0$, $\Omega$ the unique invariant vector, and Wightman fields whose
Wightman functions analytically continue to the given Schwinger functions. $\cH$ is the completion of the space of
``positive-time'' functionals $F=\sum_ic_if_i$ (finite sums of tensor products of test functions supported in
$\{x_0>0\}$) modulo null vectors, with
$\ip{\Psi_F}{e^{-tH}\Psi_G}=S(\Theta F\otimes G_t)$, where $\Theta$ is time reflection combined with complex conjugation
and $G_t$ is $G$ translated by $t$ in time. (Osterwalder--Schrader 1973, 1975; Glimm--Jaffe, Theorems 6.1.5 and 19.1.3.)
\end{citedtheorem}

\begin{theorem}[Reduction]\label{thm:reduction}
Assume Hypotheses~\ref{hyp:L1}--\ref{hyp:L3}. Then the limiting Schwinger functions satisfy (OS0)--(OS4); the theory
reconstructed by Cited Theorem~\ref{thm:OS} is non-trivial and its Hamiltonian satisfies
\[
\spec(H)\cap(0,\Delta)=\emptyset .
\]
\end{theorem}

\begin{proof}
(OS0), (OS1) are Hypothesis~\ref{hyp:L1}. (OS3): the lattice observables are commuting bounded functions, so every
$S^{(n)}$ is symmetric in its arguments (with the index $\alpha_i$ attached to $x_i$), and symmetry passes to the limit.

(OS2): Let $F=\sum_ic_if_i$ be a positive-time functional with each $f_i$ a tensor product of test functions supported
in $\{x_0>0\}$. For $n$ large the supports lie in the positive region $\{0<x_0\le T\}$ of the reflection of
Section~\ref{sec:rp} on $\Lambda_n$ (with $L^{(n)}_0$ even, which we may assume by passing to a subsequence of lattices
or doubling $L_0$; the hypotheses are insensitive to this). By reflection covariance of the observables and
antilinearity of $\Theta$, $S^{(n)}(\Theta F\otimes F)=\langle\Theta(\Phi_F)\Phi_F\rangle_{\Lambda_n,\beta_n}$ with
$\Phi_F:=\sum_ic_i\Phi^{(n)}(f_i)\in\cA_+$ the corresponding lattice functional, which is $\ge0$ by Theorem~\ref{thm:rp}.
Passing to the limit $n\to\infty$ (Hypothesis~\ref{hyp:L1}) gives $S(\Theta F\otimes F)\ge0$.

(OS4): Hypothesis~\ref{hyp:L2} passes to the limit and gives $|S(f\otimes g_t)-S(f)S(g)|\le Ce^{-\Delta t}$ for
translations in the time direction; by Euclidean invariance of the limit (OS1), the same bound holds for translations
in any direction (rotate the direction to $e_0$, which transforms the test functions but not the constants' finiteness).
In particular $S(f\otimes g_{\lambda a})\to S(f)S(g)$ as $\lambda\to\infty$.

Non-triviality is Hypothesis~\ref{hyp:L3}: a generalised free field has vanishing truncated four-point functions.

\emph{Mass gap.} In the reconstructed theory take a positive-time functional $F$ and
$\psi:=\Psi_F-\ip\Omega{\Psi_F}\Omega$. Using the semigroup formula of Cited Theorem~\ref{thm:OS} and
$\ip{\Psi_F}{\Omega}=S(\Theta F)$,
\[
\ip\psi{e^{-tH}\psi}=S(\Theta F\otimes F_t)-S(\Theta F)\,S(F)\qquad(t\ge0),
\]
and by the limiting form of Hypothesis~\ref{hyp:L2} (with $f=\Theta F$, $g=F$, both finite sums of tensor products of
test functions) the right-hand side is bounded by $C_Fe^{-\Delta t}$. The spectral argument in the proof of
Theorem~\ref{thm:gapclust}(ii) shows that the spectral measure of $H$ for $\psi$ gives no mass to $(0,\Delta)$
(the vacuum component has been removed, so no mass to $[0,\Delta)$). Since the vectors $\Psi_F$ are dense in $\cH$ by
construction and $\psi\mapsto\Psi_F$ is continuous, the vectors $\psi$ are dense in $\Omega^\perp$, and the spectral
projection of $H$ onto $(0,\Delta)$ (whose range lies in $\Omega^\perp$) vanishes.
\end{proof}

\begin{remark}[What the theorem does and does not say]
\begin{enumerate}[label=(\roman*)]
\item The gauge-invariant Hilbert space is obtained directly; there are no unphysical states to remove.
\item The hard content is entirely in Hypotheses~\ref{hyp:L1} and~\ref{hyp:L2}. Hypothesis~\ref{hyp:L1} is the
ultraviolet problem (currently: Balaban's bounds, no convergence theorem). Hypothesis~\ref{hyp:L2} contains both the
infinite-volume limit and the mass gap, uniformly along the trajectory $\beta_n\to\infty$.
\item The lattice gap $m_{\Lambda_s}(\beta)$ of Definition~\ref{def:H} is related to $\Delta$ by
$\Delta=\liminf m_{\Lambda_s^{(n)}}(\beta_n)/a_n$; Hypothesis~\ref{hyp:L2} is the statement that this is positive together
with uniform constants for the correlation functions. This is conjecture (G2) of \texttt{docs/01-problem-statement.md}.
\end{enumerate}
\end{remark}

\section{Strong coupling: a mass gap uniform in the volume}\label{sec:strong}

We prove that for $\beta$ below an explicit threshold the exponential clustering of Theorem~\ref{thm:gapclust}
holds with constants independent of the lattice, so that $\inf_{\Lambda_s}m_{\Lambda_s}(\beta)>0$. The argument is
Dobrushin's uniqueness method in a self-contained coupling form; it uses nothing about the group beyond compactness
and the bound $|{\log w_\beta}|\le\beta$, and therefore (as it must) it says nothing about large $\beta$ and does not
distinguish abelian from non-abelian groups. The constants are far from optimal.

\subsection{An abstract covariance bound}

\begin{definition}[Finite spin system]\label{def:spin}
Let $V$ be a finite set, $(S,\sigma)$ a compact metric space with a Borel probability measure, and
$\rho:S^V\to(0,\infty)$ continuous. Put $\mu:=\rho\,\sigma^{\otimes V}/\int\rho\,d\sigma^{\otimes V}$. For $x\in V$,
$\omega\in S^V$ and $s\in S$ write $\omega^{x\leftarrow s}$ for $\omega$ with the $x$-coordinate replaced by $s$, and let
\[
\mu_x(ds\,|\,\omega):=\frac{\rho(\omega^{x\leftarrow s})\,\sigma(ds)}{\int\rho(\omega^{x\leftarrow s'})\,\sigma(ds')}
\]
(the single-site conditional distribution; it depends on $\omega$ only through $\omega_{V\setminus\{x\}}$). For probability
measures $\pi,\pi'$ on $S$ let $\norm{\pi-\pi'}_\TV:=\sup_B|\pi(B)-\pi'(B)|$. The \emph{Dobrushin matrix} is
\[
C_{xy}:=\sup\bigl\{\norm{\mu_x(\cdot|\omega)-\mu_x(\cdot|\omega')}_\TV:\ \omega,\omega'\text{ agree off }y\bigr\}\ (x\ne y),
\]
$C_{xx}:=0$, and $\alpha:=\max_{x}\sum_yC_{xy}$.
For $f:S^V\to\C$ bounded measurable, $\delta_x(f):=\sup\{|f(\omega)-f(\omega')|:\omega,\omega'\text{ agree off }x\}$, and
$f$ \emph{depends only on} $A\subseteq V$ if $f(\omega)=f(\omega')$ whenever $\omega,\omega'$ agree on $A$.
\end{definition}

Two elementary facts: if $\omega,\omega'$ differ exactly on $D\subseteq V$, then, interpolating one coordinate at a time,
\begin{equation}\label{eq:interp}
|f(\omega)-f(\omega')|\le\sum_{x\in D}\delta_x(f),\qquad
\norm{\mu_x(\cdot|\omega)-\mu_x(\cdot|\omega')}_\TV\le\sum_{y\in D}C_{xy}.
\end{equation}

\begin{lemma}[Maximal coupling]\label{lem:coupling}
For probability measures $\pi,\pi'$ on $S$ with continuous densities $p,p'$ with respect to $\sigma$, there is a probability
measure $\gamma$ on $S\times S$ with marginals $\pi,\pi'$ and $\gamma(\{s\ne s'\})=\norm{\pi-\pi'}_\TV$; it can be chosen
to depend measurably on $(p,p')$.
\end{lemma}

\begin{proof}
Let $q:=\min(p,p')$, $\tau:=1-\int q\,d\sigma=\int(p-q)=\int(p'-q)=\norm{\pi-\pi'}_\TV$. If $\tau=0$ let $\gamma$ be the
diagonal image of $\pi$. Otherwise let $\gamma:=\mathrm{diag}_*(q\sigma)+\tau^{-1}\,(p-q)\sigma\otimes(p'-q)\sigma$. The marginals are
$q\sigma+(p-q)\sigma=\pi$ and $\pi'$; the off-diagonal mass is $\tau$ because $(p-q)(p'-q)=0$ pointwise, so the second term
gives no mass to the diagonal. The construction is explicit and measurable in $(p,p')$.
\end{proof}

\begin{lemma}[Ergodicity of the Gibbs sampler]\label{lem:ergodic}
Let $\emptyset\ne V'\subseteq V$ and $\eta\in S^{V\setminus V'}$, and let $\mu^\eta$ be the probability measure on $S^{V'}$ with
density $\propto\rho(\cdot,\eta)$. Let $Q$ be the Markov kernel on $S^{V'}$ of the random-scan Gibbs sampler: choose $x\in V'$
uniformly and replace $\omega_x$ by a sample from $\mu_x(\cdot|(\omega,\eta))$. Then $\mu^\eta$ is $Q$-invariant and
$\sup_\omega\norm{Q^k(\omega,\cdot)-\mu^\eta}_\TV\to0$ as $k\to\infty$.
\end{lemma}

\begin{proof}
Each single-site update is reversible with respect to $\mu^\eta$ (it resamples a coordinate from its conditional
distribution), so $\mu^\eta$ is invariant. Let $n=|V'|$ and $\epsilon:=\min\rho/\max\rho>0$; every single-site conditional
density is bounded below by $\epsilon$. With probability $n!/n^n>0$ the first $n$ site choices visit every site of $V'$;
on that event the final configuration has, conditionally on the choices, a density $\ge\epsilon^n$ with respect to
$\sigma^{\otimes V'}$ (the last update of each site produces its final value with conditional density $\ge\epsilon$ given
everything before). Hence $Q^n(\omega,\cdot)\ge c\,\sigma^{\otimes V'}$ with $c=\epsilon^nn!/n^n$ for all $\omega$. By the
standard Doeblin argument this minorisation gives $\norm{Q^{nk}(\omega,\cdot)-Q^{nk}(\omega',\cdot)}_\TV\le(1-c)^k$ uniformly in
$\omega,\omega'$, hence uniform convergence to the unique invariant measure $\mu^\eta$.
\end{proof}

\begin{theorem}[Covariance bound under Dobrushin's condition]\label{thm:dobrushin}
In the setting of Definition~\ref{def:spin} assume $\alpha<1$ and put $D:=\sum_{n\ge1}C^n$ (entrywise convergent, since
$C\ge0$ has row sums $\le\alpha$). Let $A,B\subseteq V$ be disjoint, $f$ bounded measurable depending only on $A$ and $g$
bounded measurable depending only on $B$. Then
\[
\bigl|\Cov_\mu(f,g)\bigr|\le\norm g_\infty\sum_{x\in A}\delta_x(f)\sum_{y\in B}D_{xy} .
\]
If moreover $C_{xy}=0$ unless $x\sim y$ for a graph structure on $V$ with graph distance $\varrho$, then
$\sum_{y\in B}D_{xy}\le\alpha^{\varrho(x,B)}/(1-\alpha)$, so that
\[
\bigl|\Cov_\mu(f,g)\bigr|\le\frac{\norm g_\infty}{1-\alpha}\sum_{x\in A}\delta_x(f)\,\alpha^{\varrho(x,B)}
\le\frac{\norm g_\infty\,\norm{\delta(f)}_1}{1-\alpha}\,\alpha^{\varrho(A,B)},\qquad \norm{\delta(f)}_1:=\sum_x\delta_x(f).
\]
\end{theorem}

\begin{proof}
\emph{Step 1 (conditional expectations differ little in the boundary condition).} Let $V':=V\setminus B$ and, for
$\eta,\eta'\in S^B$, let $\mu^\eta,\mu^{\eta'}$ be as in Lemma~\ref{lem:ergodic}. Run two copies of the random-scan Gibbs
sampler, $(\omega^k)$ with boundary $\eta$ and $(\omega'^k)$ with boundary $\eta'$, from a common initial configuration
$\omega^0=\omega'^0$, with the same uniformly chosen sites $x_k$, and at each step coupling the two resampled values by the
maximal coupling of Lemma~\ref{lem:coupling}. Let $n=|V'|$ and $r_k(x):=\mathbb P(\omega^k_x\ne\omega'^k_x)$ for $x\in V'$, and
set $\tilde r_k(y):=r_k(y)$ for $y\in V'$, $\tilde r_k(y):=1$ for $y\in B$. At step $k+1$, if $x_{k+1}\ne x$ the status of $x$
is unchanged; if $x_{k+1}=x$ the new values differ with conditional probability
$\norm{\mu_x(\cdot|(\omega^k,\eta))-\mu_x(\cdot|(\omega'^k,\eta'))}_\TV\le\sum_{y}C_{xy}\mathbf 1[\text{disagreement at }y]$
by~\eqref{eq:interp}. Taking expectations,
\[
r_{k+1}(x)\le\Bigl(1-\frac1n\Bigr)r_k(x)+\frac1n\sum_{y\in V}C_{xy}\,\tilde r_k(y)\qquad(x\in V').
\]
Let $r^*:=\sum_{m\ge0}(C_{V'V'})^mC_{V'B}\mathbf 1_B$ (convergent since $\norm{C_{V'V'}}_{\infty\to\infty}\le\alpha<1$); it
satisfies $r^*=C_{V'V'}r^*+C_{V'B}\mathbf 1_B$. By induction $r_k\le r^*$ for all $k$: $r_0=0\le r^*$, and if $r_k\le r^*$ then
$r_{k+1}\le(1-\frac1n)r^*+\frac1n(C_{V'V'}r^*+C_{V'B}\mathbf 1_B)=r^*$, using $C\ge0$. Moreover
$r^*_x\le\sum_{y\in B}\sum_{m\ge0}(C^{m+1})_{xy}=\sum_{y\in B}D_{xy}$. By~\eqref{eq:interp},
\[
\bigl|\E f(\omega^k)-\E f(\omega'^k)\bigr|\le\E\sum_{x\in A}\delta_x(f)\mathbf 1[\omega^k_x\ne\omega'^k_x]
=\sum_{x\in A}\delta_x(f)\,r_k(x)\le\sum_{x\in A}\delta_x(f)\sum_{y\in B}D_{xy}.
\]
By Lemma~\ref{lem:ergodic}, $\E f(\omega^k)\to\mu^\eta(f)$ and $\E f(\omega'^k)\to\mu^{\eta'}(f)$, so
$|\mu^\eta(f)-\mu^{\eta'}(f)|\le\sum_{x\in A}\delta_x(f)\sum_{y\in B}D_{xy}$ for all $\eta,\eta'$.

\emph{Step 2 (covariance).} The map $\eta\mapsto\mu^\eta(f)$ is a version of the conditional expectation $\mu(f\,|\,\omega_B)$,
and $\mu(f)=\int\mu^{\omega'_B}(f)\,\mu(d\omega')$. Since $g$ depends only on $B$,
\[
\Cov_\mu(f,g)=\int g(\omega)\bigl[\mu^{\omega_B}(f)-\mu(f)\bigr]\mu(d\omega),\qquad
\bigl|\mu^{\omega_B}(f)-\mu(f)\bigr|\le\sup_{\eta'}\bigl|\mu^{\omega_B}(f)-\mu^{\eta'}(f)\bigr|,
\]
and Step~1 gives the first bound.

\emph{Step 3 (graph distance).} If $C_{xy}=0$ unless $x\sim y$, then $(C^n)_{xy}=0$ unless $\varrho(x,y)\le n$, while
$\sum_y(C^n)_{xy}\le\alpha^n$ (row sums). Hence $\sum_{y\in B}D_{xy}\le\sum_{n\ge\max(1,\varrho(x,B))}\alpha^n\le\alpha^{\varrho(x,B)}/(1-\alpha)$,
using $\varrho(x,B)\ge1$ for $x\notin B$.
\end{proof}

\subsection{Application to the Wilson measure}

Let all $L_\mu\ge3$. Two distinct links $\ell\ne\ell'$ are \emph{adjacent}, $\ell\sim\ell'$, if some plaquette contains both;
then that plaquette is unique (two distinct unit segments lie in at most one unit square when all $L_\mu\ge3$), and each
link has exactly $K:=6(d-1)$ adjacent links (the $2(d-1)$ plaquettes through $\ell$ contribute three further links each, all
distinct). Let $\varrho$ be the graph distance on $E(\Lambda)$.

\begin{lemma}\label{lem:dobrushinbound}
For the Wilson measure (Definition~\ref{def:wilson}) viewed as a finite spin system with $V=E(\Lambda)$, $S=G$, $\sigma=dU$,
$\rho=W_\beta$: $C_{\ell\ell'}=0$ unless $\ell\sim\ell'$, and $C_{\ell\ell'}\le e^{4\beta}-1$ for $\ell\sim\ell'$. Hence
\[
\alpha\le\alpha(\beta):=6(d-1)\bigl(e^{4\beta}-1\bigr).
\]
\end{lemma}

\begin{proof}
$\mu_\ell(dU_\ell|\omega)\propto\prod_{p\ni\ell}w_\beta(U_p)\dU_\ell$ depends only on the links of the plaquettes through $\ell$, so
$C_{\ell\ell'}=0$ unless $\ell\sim\ell'$. If $\ell\sim\ell'$ and $\omega,\omega'$ agree off $\ell'$, the two conditional densities are
$e^{\phi}/Z$ and $e^{\phi'}/Z'$ with $\phi-\phi'=\log w_\beta(U_p)-\log w_\beta(U_p')$ for the unique common plaquette $p$, so
$\sup|\phi-\phi'|\le2\beta=:\delta$ by~\eqref{eq:w}. Then $Z'/Z=\int e^{\phi'-\phi}\,e^\phi/Z\in[e^{-\delta},e^\delta]$, so the ratio of
densities lies in $[e^{-2\delta},e^{2\delta}]$, and
$\norm{\mu_\ell(\cdot|\omega)-\mu_\ell(\cdot|\omega')}_\TV=\int(e^{\phi'}/Z'-e^\phi/Z)_+\le(e^{2\delta}-1)\int e^\phi/Z=e^{4\beta}-1$.
\end{proof}

\begin{theorem}[Strong-coupling clustering and mass gap, uniform in the volume]\label{thm:strong}
Let $d\ge2$, all $L_\mu\ge3$, and
\[
0<\beta<\beta_D(d):=\tfrac14\log\Bigl(1+\frac1{6(d-1)}\Bigr)\qquad(\beta_D(4)=\tfrac14\log\tfrac{19}{18}\approx0.0135),
\]
so that $\alpha(\beta)=6(d-1)(e^{4\beta}-1)<1$. Then:
\begin{enumerate}[label=(\roman*)]
\item (Clustering in every direction) For bounded measurable $f,g$ on $\Omega_\Lambda$ depending only on disjoint link sets $A,B$,
\[
\bigl|\langle fg\rangle_{\Lambda,\beta}-\langle f\rangle_{\Lambda,\beta}\langle g\rangle_{\Lambda,\beta}\bigr|
\le\frac{\norm g_\infty\norm{\delta(f)}_1}{1-\alpha(\beta)}\,\alpha(\beta)^{\varrho(A,B)} .
\]
The constant does not depend on $\Lambda$.
\item (Mass gap) For every spatial lattice $\Lambda_s$ with all $L_i\ge3$,
\[
m_{\Lambda_s}(\beta)\ \ge\ \log\frac1{\alpha(\beta)}\ >0 .
\]
In particular $\inf_{\Lambda_s}m_{\Lambda_s}(\beta)\ge\log(1/\alpha(\beta))\to\infty$ as $\beta\downarrow0$.
\end{enumerate}
\end{theorem}

\begin{proof}
(i) is Theorem~\ref{thm:dobrushin} with Lemma~\ref{lem:dobrushinbound}. (ii) For a link $\ell=(x,\mu)$ let $t(\ell):=x_0$.
The four links of a plaquette have $t$-values in $\{x_0,x_0+1\}$ for its base point $x$, so adjacent links satisfy
$|t(\ell)-t(\ell')|\le1$ modulo $L_0$, and $\varrho(\ell,\ell')\ge\min(|t(\ell)-t(\ell')|,L_0-|t(\ell)-t(\ell')|)$. Let $A$ be a
bounded gauge-invariant function on $\Omega_s$ and $0\le t\le L_0/2$. Then $\bar A(t)$ and $A(0)$ depend on the spatial links of the
slices $t$ and $0$, at graph distance $\ge t$, so by (i)
\[
\bigl|\langle\bar A(t)A(0)\rangle-\langle\bar A(t)\rangle\langle A(0)\rangle\bigr|\le\frac{\norm A_\infty\norm{\delta(A)}_1}{1-\alpha}\,\alpha^{t}
\qquad(0<t\le L_0/2),
\]
and trivially for $t=0$ with the same constant enlarged by $\norm A_\infty^2$. The constant is independent of $L_0$ and
$\norm{\delta(A)}_1\le2|E_s|\norm A_\infty<\infty$. Corollary~\ref{cor:gapfromcorr} with $m=\log(1/\alpha)$ gives the claim.
\end{proof}

\begin{remark}[On the constants and on what this proves]
\begin{enumerate}[label=(\roman*)]
\item The true strong-coupling region extends to $\beta\approx2$ for $SU(2)$ (where the Monte Carlo data in \texttt{results/}
leave the strong-coupling series), i.e.\ two orders of magnitude beyond $\beta_D$; the cluster expansion of
Section~\ref{sec:polymer} has a better but still small radius, and the physical rate is $m\approx4\log(1/u(\beta))$ with
$u(\beta)=c_\rho(\beta)/(Nc_0(\beta))\approx\beta/(2N)$, against $\log(1/\alpha)\approx\log(1/(24\beta))$ here.
\item The proof uses only $|{\log w_\beta}|\le\beta$ and compactness. It applies verbatim to $U(1)$, where the conclusion fails
at large $\beta$ (Guth; Fröhlich--Spencer). Thus no refinement of this argument can be the mechanism for the mass gap at
large $\beta$; it establishes the strong-coupling endpoint of the problem and nothing more.
\item Statement (i) also yields, with a little more work (comparison of different volumes through the same coupling
argument, with the boundary playing the role of $B$), the existence of the infinite-volume limit of all local expectations
for $\beta<\beta_D$; this is milestone T1.1 of the research program and is not carried out here.
\end{enumerate}
\end{remark}

\section{The polymer representation and the cluster expansion}\label{sec:polymer}

The Dobrushin argument of Section~\ref{sec:strong} is the shortest route to a uniform gap, but it hides the
geometric structure of the strong-coupling regime. That structure---the Wilson measure as a gas of closed
surfaces---is what the physical pictures of confinement (flux tubes, glueballs as closed tubes) are built on,
and it is the natural endpoint for any renormalisation-group argument. We derive it and prove convergence of
the associated cluster expansion for small $\beta$, using the Kotecký--Preiss criterion.

\subsection{Polymers}

\begin{definition}\label{def:polymer}
Let $f_\beta:=w_\beta/c_0(\beta)-1$, so that $w_\beta=c_0(\beta)(1+f_\beta)$, $\int f_\beta\dU=0$, and, by
\eqref{eq:w} and Corollary~\ref{cor:uniform} ($1\le c_0\le e^\beta$),
\begin{equation}\label{eq:eps}
\norm{f_\beta}_\infty\le\epsilon(\beta):=e^{2\beta}-1 .
\end{equation}
Two plaquettes are \emph{adjacent} if they share a link. A \emph{polymer} is a non-empty set $\gamma\subseteq\cP(\Lambda)$
which is connected under adjacency; $|\gamma|$ is its number of plaquettes and $E(\gamma)$ the set of links of its
plaquettes. Two polymers are \emph{compatible}, $\gamma\approx\gamma'$, if $E(\gamma)\cap E(\gamma')=\emptyset$; otherwise
$\gamma\not\approx\gamma'$ (in particular $\gamma\not\approx\gamma$). The \emph{activity} of $\gamma$ is
\[
z(\gamma):=\int\prod_{p\in\gamma}f_\beta(U_p)\dU\qquad(\text{an integral over }G^{E(\gamma)}).
\]
\end{definition}

\begin{proposition}[Polymer representation]\label{prop:polymer}
\[
Z_\Lambda(\beta)=c_0(\beta)^{|\cP(\Lambda)|}\,\Xi_\Lambda(\beta),\qquad
\Xi_\Lambda(\beta):=\sum_{n\ge0}\frac1{n!}\sum_{\substack{\gamma_1,\dots,\gamma_n\\ \gamma_i\approx\gamma_j\ (i\ne j)}}\prod_{i=1}^nz(\gamma_i),
\]
the $n=0$ term being $1$ (the sum is over ordered $n$-tuples of pairwise compatible polymers).
\end{proposition}

\begin{proof}
Expanding $\prod_p(1+f_\beta(U_p))$, $Z_\Lambda/c_0^{|\cP|}=\sum_{S\subseteq\cP}\int\prod_{p\in S}f_\beta(U_p)\dU$. Each $S$
decomposes uniquely into its connected components $\gamma_1,\dots,\gamma_n$ (connected under adjacency), which are
pairwise compatible (two components sharing a link would be adjacent, hence one component). Distinct components
depend on disjoint sets of link variables, so the integral factorises into $\prod_iz(\gamma_i)$. Conversely, a family of
pairwise compatible polymers is the set of components of its union. Counting ordered tuples introduces the factor $1/n!$.
\end{proof}

\begin{lemma}[Closed surfaces]\label{lem:closed}
$z(\gamma)=0$ unless every link $\ell\in E(\gamma)$ belongs to at least two plaquettes of $\gamma$. For every polymer,
$|z(\gamma)|\le\epsilon(\beta)^{|\gamma|}$.
\end{lemma}

\begin{proof}
If $\ell$ lies in exactly one plaquette $p\in\gamma$, then $U_p=AU_\ell^{\pm1}B$ with $A,B$ independent of $U_\ell$, and by
left, right and inversion invariance of Haar measure $\int f_\beta(AU_\ell^{\pm1}B)\dU_\ell=\int f_\beta\dU=0$; by Fubini,
$z(\gamma)=0$. The bound is \eqref{eq:eps}.
\end{proof}

\begin{lemma}[Counting polymers]\label{lem:count}
Let $K_\cP:=4(2d-3)$, the maximal number of plaquettes adjacent to a given plaquette. For every plaquette $p_0$ and
$n\ge1$, the number of polymers $\gamma$ with $p_0\in\gamma$ and $|\gamma|=n$ is at most $(eK_\cP)^n$.
\end{lemma}

\begin{proof}
A plaquette has four links, each contained in $2(d-1)$ plaquettes, so at most $4(2(d-1)-1)=K_\cP$ other plaquettes are
adjacent to it. Fix, for every plaquette, an enumeration of its adjacent plaquettes by $\{1,\dots,K_\cP\}$. Given $\gamma\ni p_0$
with $|\gamma|=n$, let $\tau$ be the spanning tree of $\gamma$ produced by depth-first search from $p_0$ that always proceeds
to the unvisited adjacent plaquette of smallest index; $\tau$ is a rooted tree in which every vertex $p$ has at most $K_\cP$
children, each child labelled by its index among the neighbours of $p$. Record $\tau$ by listing its vertices in depth-first
order and, for each vertex, the $K_\cP$-bit string indicating which indices are occupied by children. The record determines
$\tau$ and hence $\gamma$ (walk from $p_0$ following the indices), so the map $\gamma\mapsto$ record is injective. A record is a
string of $K_\cP n$ bits with exactly $n-1$ ones, so the number of polymers is at most
$\binom{K_\cP n}{n-1}\le\binom{K_\cP n}{n}\le(eK_\cP)^n$ (binomial coefficients $\binom mk$ increase in $k$ for
$k\le(m+1)/2$, and $n\le(K_\cP n+1)/2$ since $K_\cP\ge2$; finally $\binom mk\le(em/k)^k$).
\end{proof}

\subsection{The Kotecký--Preiss criterion}

\begin{citedtheorem}[Kotecký--Preiss 1986]\label{thm:KP}
Let $\mathcal P$ be a finite set, $\not\approx$ a symmetric reflexive relation on it (``incompatibility''), $z:\mathcal P\to\C$,
and $\Xi:=\sum_{n\ge0}\frac1{n!}\sum_{\gamma_1,\dots,\gamma_n\text{ pairwise compatible}}\prod_iz(\gamma_i)$. Suppose
$a:\mathcal P\to[0,\infty)$ satisfies
\begin{equation}\label{eq:KPcond}
\sum_{\gamma'\not\approx\gamma}|z(\gamma')|\,e^{a(\gamma')}\le a(\gamma)\qquad\text{for all }\gamma\in\mathcal P .
\end{equation}
Then $\Xi\neq0$ and
\[
\log\Xi=\sum_{n\ge1}\frac1{n!}\sum_{\gamma_1,\dots,\gamma_n}\varphi^T(\gamma_1,\dots,\gamma_n)\prod_{i=1}^nz(\gamma_i),
\]
where $\varphi^T(\gamma_1,\dots,\gamma_n)=\sum_{H}(-1)^{|H|}$ is the Ursell function (sum over connected spanning subgraphs
$H$ of the graph on $\{1,\dots,n\}$ with an edge $\{i,j\}$ whenever $\gamma_i\not\approx\gamma_j$; $\varphi^T=1$ for $n=1$), the
series converges absolutely, and for every $\gamma\in\mathcal P$
\[
\sum_{n\ge1}\frac1{n!}\sum_{\substack{\gamma_1,\dots,\gamma_n\\ \exists i:\ \gamma_i\not\approx\gamma}}
|\varphi^T(\gamma_1,\dots,\gamma_n)|\prod_i|z(\gamma_i)|\le a(\gamma).
\]
(Kotecký--Preiss, Commun.\ Math.\ Phys.\ 103 (1986); see also Friedli--Velenik, \emph{Statistical Mechanics of Lattice
Systems}, Theorem~5.4, and Ueltschi, Moscow Math.\ J.\ 4 (2004).)
\end{citedtheorem}

\begin{theorem}[Convergence of the strong-coupling cluster expansion]\label{thm:conv}
Let $d\ge2$ and
\[
0<\beta<\beta_0(d):=\tfrac12\log\Bigl(1+\frac1{e^2K_\cP(K_\cP+2)}\Bigr),\qquad K_\cP=4(2d-3).
\]
Then \eqref{eq:KPcond} holds for the polymer system of Definition~\ref{def:polymer} with $a(\gamma)=|\gamma|$, on every
finite lattice $\Lambda$. Consequently $\log\Xi_\Lambda(\beta)$ is given by the absolutely convergent cluster series of
Cited Theorem~\ref{thm:KP}, and for every plaquette $p_0$ the total weight of clusters touching $p_0$ is at most~$1$:
\[
\sum_{n\ge1}\frac1{n!}\sum_{\substack{\gamma_1,\dots,\gamma_n\\ \exists i:\ p_0\in\gamma_i\text{ or }p_0\text{ adjacent to }\gamma_i}}
|\varphi^T(\gamma_1,\dots,\gamma_n)|\prod_i|z(\gamma_i)|\le1 .
\]
\end{theorem}

\begin{proof}
Let $x:=e^2K_\cP\,\epsilon(\beta)$; the hypothesis $\beta<\beta_0(d)$ is exactly $x<1/(K_\cP+2)$, so that
$(K_\cP+1)\,x/(1-x)\le1$. Fix a polymer $\gamma_0$. A polymer $\gamma\not\approx\gamma_0$ shares a link with $\gamma_0$, hence contains
a plaquette $p$ that is equal or adjacent to a plaquette of $\gamma_0$; the set $N(\gamma_0)$ of such $p$ has at most
$(K_\cP+1)|\gamma_0|$ elements. By Lemmas~\ref{lem:closed} and~\ref{lem:count},
\[
\sum_{\gamma\not\approx\gamma_0}|z(\gamma)|e^{|\gamma|}
\le\sum_{p\in N(\gamma_0)}\sum_{n\ge1}\sum_{\substack{\gamma\ni p\\ |\gamma|=n}}\epsilon^ne^n
\le(K_\cP+1)|\gamma_0|\sum_{n\ge1}(eK_\cP)^n\epsilon^ne^n=(K_\cP+1)|\gamma_0|\frac{x}{1-x}\le|\gamma_0|=a(\gamma_0).
\]
The last statement is the final bound of Cited Theorem~\ref{thm:KP} applied to $\gamma=\{p_0\}$ (a set of one plaquette,
which we may add to the polymer system with activity $0$ without affecting anything: the clusters incompatible with it are
exactly those touching $p_0$, and $a(\{p_0\})=1$).
\end{proof}

\begin{corollary}[Free energy]\label{cor:free}
For $0<\beta<\beta_0(d)$ the limit
\[
F(\beta):=\lim_{L_0,\dots,L_{d-1}\to\infty}\frac{1}{|\cP(\Lambda)|}\log Z_\Lambda(\beta)
=\log c_0(\beta)+\sum_{n\ge1}\frac1{n!}\sum_{\substack{\gamma_1,\dots,\gamma_n\subset\cP(\Z^d)\\ p_0\in\gamma_1}}
\frac{\varphi^T(\gamma_1,\dots,\gamma_n)}{|\gamma_1\cup\cdots\cup\gamma_n|}\prod_iz(\gamma_i)
\]
exists and is given by the absolutely convergent series on the right (independent of the plaquette $p_0$ of $\Z^d$).
\end{corollary}

\begin{proof}
By Proposition~\ref{prop:polymer} and Theorem~\ref{thm:conv}, $\frac1{|\cP|}\log Z_\Lambda=\log c_0+\frac1{|\cP|}\sum_X\Psi(X)$,
where $X$ runs over clusters (tuples) on $\Lambda$ and $\Psi(X)=\frac1{n!}\varphi^T(X)\prod z(\gamma_i)$. Writing
$\frac1{|\cP|}\sum_X\Psi(X)=\frac1{|\cP|}\sum_{p\in\cP}\sum_{X\ni p}\Psi(X)/|\supp X|$ (where $\supp X=\bigcup\gamma_i$ and $X\ni p$
means $p\in\supp X$) and using translation invariance, this equals $\sum_{X\ni p_0}\Psi(X)/|\supp X|$ on the torus. A cluster
on the torus whose support has diameter less than $\min_\mu L_\mu$ corresponds to exactly one cluster in $\Z^d$ with the same
activities (the activity depends only on the combinatorial structure of the polymers), and the clusters of diameter
$\ge\min_\mu L_\mu$ have total weight at most $\sum_{m\ge\min L_\mu}(\text{weight of clusters touching }p_0\text{ of size }\ge m)$,
which tends to $0$ as $L\to\infty$ by absolute convergence (the size of a cluster is at least its diameter). The series in $\Z^d$
converges absolutely by the same bound as in Theorem~\ref{thm:conv}, which does not depend on the lattice size.
\end{proof}

\subsection{The first polymer, and a check against the simulations}

\begin{proposition}[Cube activity]\label{prop:cube}
Let $d\ge3$ and let $\gamma_{\rm cube}$ be the six faces of a unit cube in $\Lambda$ (with all $L_\mu\ge3$). Then
\[
z(\gamma_{\rm cube})=\sum_{r\neq0}\Bigl(\frac{c_r(\beta)}{d_r\,c_0(\beta)}\Bigr)^{6}d_r^{\,2}.
\]
For $G=SU(2)$ this is $4u(\beta)^6+9u_1(\beta)^6+\cdots$ with $u=I_2(\beta)/I_1(\beta)$, $u_1=I_3(\beta)/I_1(\beta)$;
for $G=SU(3)$ it is $18u(\beta)^6+\cdots$ with $u=c_{\rm fund}/(3c_0)$ (fundamental and anti-fundamental contribute equally).
\end{proposition}

\begin{proof}
Expand each $f_\beta(U_p)=\sum_{r\ne0}(c_r/c_0)\chi_r(U_p)$ (Corollary~\ref{cor:uniform}; the sums may be interchanged
with the integral by uniform convergence) and orient the six faces as the boundary of the cube, so that each edge is
traversed in opposite directions by the two faces containing it; $\chi_r(U_p)$ is unchanged by cyclic relabelling of the
base point and $\chi_r(U_p^{-1})=\overline{\chi_r(U_p)}$ corresponds to the opposite orientation, so we may use the boundary
orientation for every face. Call the faces $B$ (bottom), $T$ (top) and $S_1,\dots,S_4$ (sides, in cyclic order), and let
$r_B,r_T,r_1,\dots,r_4$ be the representations attached to them. Integrate the vertical edges in the order
$S_1|S_2$, $S_2|S_3$, $S_3|S_4$: by Theorem~\ref{thm:pw}(iv) each integration vanishes unless the two adjacent side
representations agree and then merges the two characters into the character of the boundary loop of the union, with a
factor $1/d_r$. After these three integrations we have $r_1=\dots=r_4=:r$ and the single character
$\chi_r(XUYU^{-1})/d_r^3$, where $U$ is the remaining vertical edge $S_4|S_1$, traversed once in each direction, and $X,Y$ are the
products of the four top edges and of the four bottom edges (in the appropriate orientations). By
Theorem~\ref{thm:pw}(v) the integral over $U$ gives $\chi_r(X)\chi_r(Y)/d_r^4$. The top face contributes $\chi_{r_T}(X^{-1})$
(opposite orientation), and $\int\chi_r(X)\chi_{r_T}(X^{-1})$ over the four top edges equals $\int\chi_r(X)\overline{\chi_{r_T}(X)}\,dX=\delta_{r r_T}$
since $X$ is Haar distributed when any one of its four edges is; likewise the bottom face forces $r_B=r$ and contributes $1$.
Altogether $\int\prod_{p}\chi_{r_p}(U_p)\dU=\delta_{\text{all }r_p=r}\,d_r^{-4}$, and multiplying by the six coefficients
$c_r/c_0$ gives $z(\gamma_{\rm cube})=\sum_{r\ne0}(c_r/c_0)^6d_r^{-4}=\sum_{r\ne0}(c_r/(d_rc_0))^6d_r^{2}$.
For $SU(2)$, $c_j/(d_jc_0)=I_{2j+1}(\beta)/I_1(\beta)$; for $SU(3)$ the fundamental and its conjugate both have
$c_r/(3c_0)=u$ and $d_r^2=9$.
\end{proof}

\begin{remark}[Plaquette expectation; numerical check]\label{rem:plaq}
Differentiating Corollary~\ref{cor:free} and keeping only the one-polymer clusters of size six (there are $\binom d3$
cubes per site and $\binom d2$ plaquettes per site, i.e.\ $(d-2)/3$ cubes per plaquette),
\[
\langle P\rangle_\beta=\frac{d}{d\beta}F(\beta)=u(\beta)+\frac{d-2}{3}\,\frac{d}{d\beta}z(\gamma_{\rm cube})+\cdots
= u+\frac{d-2}{3}\cdot6z_6\,u^5u'(\beta)+O(u^9),
\]
with $z_6=4$ for $SU(2)$ and $z_6=18$ for $SU(3)$, i.e.\ $\langle P\rangle=u+16u^5u'$ for $SU(2)$ and $u+72u^5u'$ for
$SU(3)$ in four dimensions, where $u(\beta)=c_0'(\beta)/c_0(\beta)$ is the exact one-plaquette expectation. The Monte Carlo
scans in \texttt{results/} reproduce these formulas: for $SU(2)$ at $\beta=1$ the simulation gives
$\langle P\rangle=0.2434(2)$ against $u=0.2402$ and $u+16u^5u'=0.2430$; for $SU(3)$ at $\beta=2$ it gives $0.1288(1)$ against
$u=0.1286$ and $u+72u^5u'=0.1288$ (see \texttt{results/plaquette\_su2\_8x4.csv}, \texttt{results/plaquette\_su3\_8x4.csv}).
The expansion is formal beyond its radius of convergence, but the agreement at these couplings is the expected check of
the polymer bookkeeping.
\end{remark}

\begin{remark}[Further consequences, not proved here]\label{rem:further}
With the same machinery one obtains, for $\beta<\beta_0(d)$: exponential clustering of all local gauge-invariant
observables with constants uniform in $\Lambda$ (an alternative to Theorem~\ref{thm:strong}, with the same structure
of proof once the expectation of a local observable is written as a ratio of polymer partition functions); the existence
and uniqueness of the infinite-volume limit; analyticity of $F$ in $\beta$ in a complex neighbourhood of $(0,\beta_0)$; and
the area law for Wilson loops in representations with non-trivial $N$-ality (the surfaces carrying the character
expansion must span the loop). These are Osterwalder--Seiler's results and are milestones T1.1 of the research program.
\end{remark}

\section{Summary, and what is open}\label{sec:open}

\subsection{What has been established}

\begin{center}
\begin{tabular}{@{}lll@{}}
\toprule
Statement & Where & Status\\
\midrule
Wilson measure, gauge invariance, Elitzur-type vanishing & Thm~\ref{thm:basic} & Theorem\\
Positivity of all character coefficients & Lem~\ref{lem:cr} & Theorem\\
Reflection positivity (site reflections, periodic time) & Thm~\ref{thm:rp} & Theorem\\
Transfer matrix: positive, trace class, strictly positive on $\cH_{\rm phys}$ & Lem~\ref{lem:K} & Theorem\\
$Z=\Tr\hT^{L_0}$; correlations as traces & Thm~\ref{thm:tm} & Theorem\\
Unique positive vacuum (Perron--Frobenius) & Thm~\ref{thm:vacuum} & Theorem\\
Hamiltonian, finite-lattice gap $m_{\Lambda_s}(\beta)>0$ & Def~\ref{def:H} & Theorem\\
Ground-state expectations as $L_0\to\infty$ & Thm~\ref{thm:gs} & Theorem\\
Gap $\Leftrightarrow$ exponential clustering in time & Thm~\ref{thm:gapclust} & Theorem\\
Gap $=$ spectral gap of the ground-state chain & Thm~\ref{thm:chain} & Theorem\\
Lattice hypotheses $\Rightarrow$ OS axioms $+$ mass gap & Thm~\ref{thm:reduction} & Theorem (modulo OS reconstruction)\\
Uniform gap for $\beta<\beta_D(d)$ & Thm~\ref{thm:strong} & Theorem\\
Polymer representation, closed surfaces & Prop~\ref{prop:polymer}, Lem~\ref{lem:closed} & Theorem\\
Convergent cluster expansion for $\beta<\beta_0(d)$ & Thm~\ref{thm:conv} & Theorem (modulo Kotecký--Preiss)\\
\bottomrule
\end{tabular}
\end{center}

\subsection{What is open}

Fix $G=SU(N)$, $d=4$, all lattices with $L_\mu\ge3$.

\begin{conjecture}[G1: gap at every coupling]
For every $\beta\in(0,\infty)$, $\ \inf_{\Lambda_s}m_{\Lambda_s}(\beta)>0$.
\end{conjecture}

\begin{conjecture}[G2: the gap survives the continuum limit]
With $a\Lambda_L(\beta)=(b_0g^2)^{-b_1/(2b_0^2)}e^{-1/(2b_0g^2)}$, $g^2=2N/\beta$, $b_0=11N/(48\pi^2)$,
$b_1=34N^2/(3(16\pi^2)^2)$:
\[
\liminf_{\beta\to\infty}\ \frac{\inf_{\Lambda_s}m_{\Lambda_s}(\beta)}{a\Lambda_L(\beta)}>0 .
\]
\end{conjecture}

\begin{conjecture}[UV: convergence]
Hypothesis~\ref{hyp:L1} holds for the renormalised plaquette field and for Wilson loops along the trajectory
$\beta=\beta(a)$ determined by asymptotic scaling.
\end{conjecture}

Theorem~\ref{thm:reduction} shows that (G2) (in the uniform form of Hypothesis~\ref{hyp:L2}) together with (UV) and
non-triviality is the Clay statement. Theorem~\ref{thm:strong} proves (G1) for $\beta<\beta_D$ by a method that also
applies to $U(1)$, for which (G1) is false at large $\beta$; the mechanism for large $\beta$ must therefore be of a
different nature. Theorem~\ref{thm:chain} shows that (G1) is equivalent to a uniform spectral gap for a reversible,
locally structured Markov chain whose stationary measure is the vacuum measure $\Omega^2dU$ on the three-dimensional
spatial lattice. The research program (\texttt{docs/03-research-program.md}) is organised around these three statements.

\end{document}
